\documentclass[12pt, oneside]{book}

\usepackage[utf8]{inputenc}
\usepackage[T2A,T1]{fontenc}
\usepackage[english,russian]{babel}
\usepackage[a4paper, left=4cm, top=2cm, right=2cm, bottom=2cm]{geometry}

\usepackage{graphicx}
\usepackage{amsmath}
\usepackage{amsfonts}
\usepackage{amssymb}
\usepackage{hyperref}
\usepackage{ragged2e}
\usepackage{setspace}
\usepackage[table]{xcolor}
\usepackage{tikz}
\usepackage{multirow}
\usepackage{float}
\usepackage{longtable}
\usetikzlibrary{arrows.meta, positioning, fit, backgrounds, calc}
\definecolor{thbg}{HTML}{F3F4F6}
\definecolor{tdborder}{HTML}{E5E7EB}

\definecolor{intbg}{HTML}{F8FAFC}
\definecolor{intborder}{HTML}{94A3B8}
\definecolor{appbg}{HTML}{FEF9F0}
\definecolor{appborder}{HTML}{D97706}
\definecolor{applabel}{HTML}{92400E}
\definecolor{compbg}{HTML}{F0FDF4}
\definecolor{compborder}{HTML}{16A34A}
\definecolor{complabel}{HTML}{166534}
\definecolor{greenhi}{HTML}{DCFCE7}
\definecolor{accentblue}{HTML}{4338CA}
\definecolor{mutedtext}{HTML}{6B7280}

\onehalfspacing
\setlength{\parindent}{1.25cm}
\everymath{\displaystyle}
\usepackage{tocloft}
\cftsetindents{section}{1em}{2em}

\renewcommand\cfttoctitlefont{\hfill\Large\bfseries}
\renewcommand\cftaftertoctitle{\hfill\mbox{}}

\setlength{\cftbeforesecskip}{0.1cm}
\setlength{\cftbeforesubsecskip}{0.1cm}
\setlength{\cftbeforetoctitleskip}{0cm}
\setlength{\cftaftertoctitleskip}{0.4cm}
\setcounter{tocdepth}{2}


\makeatletter
\renewcommand{\@makeschapterhead}[1]{%
        \vspace*{-\parskip}
        \vspace*{-\baselineskip}
        \vspace*{-0.5cm}
        {\parindent \z@ \center
                \normalfont
                \LARGE \bfseries #1\par\nobreak
                \addcontentsline{toc}{chapter}{#1}
                \vskip 10\p@
}}
\makeatother

\makeatletter
\renewcommand{\chapter}{%
        \if@openright\cleardoublepage\else\clearpage\fi
        \thispagestyle{plain}
        \global\@topnum\z@
        \@afterindentfalse
\secdef\@chapter\@schapter}

\renewcommand{\@makechapterhead}[1]{%
        \vspace*{0\p@}
        {\parindent \z@ \raggedright
                \normalfont
                \LARGE \bfseries \thechapter\quad #1\par\nobreak
                \vskip 20\p@
}}
\makeatother

\makeatletter
\newcommand{\custompagestyle}{%
        \begingroup
        \renewcommand{\ps@plain}{%
                \renewcommand{\@oddfoot}{\hfil\thepage\hfil}
                \renewcommand{\@evenfoot}{\hfil\thepage\hfil}
                \renewcommand{\@oddhead}{}
                \renewcommand{\@evenhead}{}
        }%
        \pagestyle{plain}
}
\makeatother

\newenvironment{customquote}
{\list{}{\leftmargin=0pt
                \rightmargin=0pt
                \topsep=0pt
                \partopsep=0pt
                \parsep=2pt
                \itemsep=0pt
        }%
        \normalfont\relax
\item\relax}
{\endlist}


\begin{document}
\thispagestyle{empty}
\begin{center}
        \resizebox{\textwidth}{!}{\textbf{САНКТ-ПЕТЕРБУРГСКИЙ ГОСУДАРСТВЕННЫЙ УНИВЕРСИТЕТ}}

        \vspace{0.5em}
        \text{\large Факультет прикладной математики -- процессов управления}
        
        \vspace{0.6em}
        \text{\large Кафедра информационных систем}
        
        

        \vspace{7em}
        \textbf{\large Маслакова Анастасия Андреевна}

        \vspace{2em}
        \textbf{\Large Выпускная квалификационная работа}
        
        \vspace{2.5em}
        \textbf{\Large «Веб-платформа интерактивной визуализации и анализа многомерных численных алгоритмов»}

        \vspace{4em}
        Уровень образования: бакалавриат \\
        Направление: 01.03.02 Прикладная математика и информатика \\
        Основная образовательная программа СВ.5005.2022:
        «Прикладная математика, фундаментальная информатика и программирование»

        \vspace{4em}
        \begin{flushright}
                Научный руководитель \\
                доктор физ.-мат. наук, профессор \\
                Матросов Александр Васильевич
        \end{flushright}

        \begin{flushright}
                Рецензент \\
                Доцент кафедры моделирования экономических   систем, \\
                кандидат физ.-мат. наук \\
                Ковшов Александр Михайлович
        \end{flushright}

        \vfill
        Санкт-Петербург\\ \the\year
\end{center}

\newpage
\custompagestyle
\setcounter{page}{2}
\renewcommand{\contentsname}{Содержание}
\tableofcontents

\newpage
\chapter*{Введение}
        \quad В современной вычислительной науке численные методы занимают центральное место. Методы линейной алгебры, разложение матриц и итерационные алгоритмы оптимизации составляют фундамент, на котором строятся решения в инженерии, физике, экономике, науке о данных и искусственном интеллекте. Глубокое понимание этих методов необходимо для специалиста, работающего с реальными вычислительными задачами. Именно численные алгоритмы лежат в основе широкого класса прикладных проблем: от решения систем линейных уравнений высокой размерности и вычисления собственных значений матриц до спектрального анализа многомерных сигналов и численного интегрирования дифференциальных уравнений.

        \quad Однако имеется существенный разрыв между потребностями специалистов и доступными инструментами. Традиционные среды для численного анализа, такие как: MATLAB, Scilab, Mathematica, обладают широкими возможностями, но их использование сопряжено с рядом ограничений. Приобретение коммерческих лицензий требует значительных финансовых затрат, локальная установка и настройка среды создают дополнительный барьер для входа, а сами инструменты, как правило, ориентированы на опытного пользователя и не предоставляют средств для пошаговой интерактивной визуализации вычислительного процесса. Это делает их малопригодными как для первичного освоения алгоритмов, так и для быстрого исследовательского прототипирования в условиях ограниченных ресурсов.

        \quad Существующие веб-ориентированные альтернативы, например онлайн - калькуляторы, вычислительные блокноты или образовательные платформы, лишь частично закрывают эту потребность. Вычислительные среды требуют знания языков программирования и настройки серверной инфраструктуры, онлайн - калькуляторы не обеспечивают работу с данными нетривиальной размерности, а большинство образовательных ресурсов ограничиваются статическими иллюстрациями и не позволяют пользователю взаимодействовать с алгоритмом напрямую. 

        \quad Технология WebAssembly (Wasm) открывает принципиально новые возможности для заполнения этой ниши. Стандартизированный бинарный формат, поддерживаемый всеми современными браузерами, позволяет компилировать высокопроизводительные библиотеки численных вычислений в код, исполняемый непосредственно на стороне клиента с производительностью, сопоставимой с нативной. Это принципиально отличает WebAssembly от предшествующих попыток перенести вычисления в браузер. В сочетании с современными средствами веб-визуализации -WebGL, D3.js, Canvas API - данная технология создаёт техническую основу для реализации платформы нового класса.
        
        \quad Настоящая дипломная работа посвящена проектированию и разработке веб-платформы интерактивной визуализации и анализа многомерных численных методов. Платформа предоставляет пользователю единую браузерную среду, в которой становится возможным не только выполнять вычисления над многомерными данными с использованием проверенных численных библиотек, но и наблюдать за поведением алгоритмов в интерактивном режиме: визуализировать промежуточные состояния, исследовать зависимость результата от входных параметров, сравнивать методы по точности и вычислительной сложности. Подобный подход позволяет совместить доступность браузерного решения с вычислительной мощностью нативных библиотек и наглядностью, необходимой для глубокого понимания алгоритмической сути рассматриваемых методов.

\chapter*{Постановка задачи}
        \quad Целью данной работы является создание прототипа платформы для интерактивного анализа и визуализации многомерных численных методов – полезного инструмента для образования, который сделает сложные алгоритмы наглядными и доступными для изучения. 

        \quad Для достижения поставленной цели необходимо выполнить следующие задачи:
        \begin{enumerate}
                \item Провести анализ существующих инструментов и платформ для численного анализа и выработать на его основе требования к разрабатываемой платформе.
                \item Исследовать технологии, применимые для реализации высокопроизводительных вычислений в браузере и обосновать выбор технического стека.
                \item Спроектировать архитектуру и реализовать прототип платформы с поддержкой интерактивной визуализации выбранных численных методов и возможностью работы с многомерными входными данными.
                \item Разработать методику оценки производительности платформы и провести вычислительный эксперимент в разных браузерах для сравнения эффективности созданной платформы.
                \item Проанализировать полученные результаты и сформулировать рекомендации для дальнейшей доработки. 
        \end{enumerate}

\chapter*{Обзор литературы}

        \quad Формирование современных подходов к созданию высокопроизводительных вычислительных веб-приложений опирается на ряд исследований, посвящённых проектированию и стандартизации технологии WebAssembly, анализу её производительности, разработке линейно-алгебраических библиотек для веба, а также реализации интерактивной научной визуализации на стороне клиента. Анализ существующих публикаций позволяет обосновать выбор технологий и архитектурных решений, принятых в настоящей работе, и определить место разрабатываемой платформы среди существующих решений.

        \quad Фундаментальной работой, определившей развитие технологии\\ WebAssembly, является статья [1], опубликованная на конференции PLDI 2017 и подготовленная инженерами из Google, Mozilla, Microsoft и Apple. В работе показано, что JavaScript, несмотря на оптимизации JIT-компиляторов, демонстрирует непредсказуемую производительность и не является надёжной целевой платформой для компиляции. WebAssembly решает эту проблему, предоставляя компактный бинарный формат с детерминированным поведением и скоростью выполнения, близкой к нативной. 

        \quad Особый интерес представляют исследования, в который были проведены эксперименты по реализации матричных вычислений в WebAssembly и их сравнению с JavaScript-реализациями. Например, в отчете [2] реализовали библиотеку для работы с векторами и матрицами на Rust, скомпилированную в WASM через wasm-pack, и показали, что на операции умножения матриц размером до 1000×1000 WASM-реализация превосходит чистый JavaScript до 10 раз в Chrome и до 50 раз в Firefox. Кроме того, в работе [3] провели более детальный анализ на семи конфигурациях реализации умножения матриц (от 256×256 до 1024×1024), включая WebAssembly с SIMD-инструкциями и гибридные конфигурации с WebGPU, и выявили, что базовый WebAssembly без оптимизаций может быть даже медленнее JavaScript (ускорение 0.41x), поскольку JIT-компиляторы выполняют агрессивную оптимизацию паттернов доступа к памяти, - только комбинация SIMD-векторизации и блочного доступа к памяти позволяет WebAssembly превзойти JavaScript на 64\%. 

        \quad Ряд публикаций подтверждает техническую жизнеспособность компиляции существующих C/C++ библиотек в WebAssembly с использованием Emscripten для реализации интерактивной научной визуализации непосредственно в браузере. Так, в работе [4] скомпилировали libtiff и LAStools в WASM и продемонстрировали интерактивную визуализацию изоповерхности набора данных объёмом около 1 ТБ полностью в браузере с применением WebGPU: поверхность из 137,5 миллионов треугольников вычисляется за 526 мс и рендерится со скоростью 30 FPS. А в [5] описан подход к кроссплатформенной визуализации научных симуляций через WebAssembly и WebGL без внешних зависимостей, где существующий C-код компилируется через Emscripten с поддержкой трёх режимов работы: нативного (OpenGL), гибридного (сервер + браузер) и полностью браузерного.

        \quad Опыт создания полноценных образовательных вычислительных сред в браузере представлен в статье [6]. Авторы реализовали браузерную среду разработки на Python, использующую Pyodide - порт CPython, скомпилированный в WebAssembly. Производительность системы составляет 3–5x от нативного CPython, что является приемлемым для образовательных задач. Практическое внедрение в школах Швейцарии подтверждает, что образовательные платформы на основе WebAssembly могут быть не только технически реализуемы , но и востребованы в учебном процессе. 

        \quad Актуальность настоящего исследования подтверждается тем, что большая часть существующих работ либо фокусируется на отдельных аспектах, например, производительности или компиляции, либо ограничена специфическими кейсами без образовательной или визуализационной составляющей. Между тем, систематического сопоставления различных подходов в контексте создания комплексной вычислительной платформы с интерактивной визуализацией и образовательной направленностью в литературе крайне мало. 

        \quad Рассмотренная литература подчёркивает как высокую значимость темы высокопроизводительных вычислений в браузере на основе WebAssembly, так и ограниченность имеющихся работ, объединяющих численные методы, интерактивную визуализацию и образовательную функциональность в комплексной, системной форме.


\chapter{Технология WebAssembly и высокопроизводительные вычисления в браузерной среде}

\section{Определение и назначение WebAssembly} 

    \quad WebAssembly (Wasm) представляет собой открытый стандарт, определяющий переносимый бинарный формат инструкций для стековой виртуальной машины. Стандарт разработан рабочей группой W3C WebAssembly Working Group при участии инженеров компаний Google, Mozilla, Microsoft и Apple и опубликован в качестве рекомендации W3C в декабре 2019 года [1]. WebAssembly проектировался как целевая платформа компиляции для языков с явным управлением памятью, таких как C, C++ и Rust, и предназначен для выполнения вычислительно интенсивных задач в браузерной среде с производительностью, приближающейся к нативной.

    \quad Необходимость появления WebAssembly связана с эволюцией требований к веб-приложениям, которые всё чаще включают ресурсоёмкие вычисления: обработку графики, моделирование, анализ данных и выполнение алгоритмов, характерных для настольного программного обеспечения. Хотя современные JavaScript-движки, такие как V8, SpiderMonkey и JavaScriptCore, используют сложные механизмы оптимизации, включая JIT-компиляцию, специфика языка накладывает ограничения на достижимую эффективность и предсказуемость выполнения. В задачах, требующих строгого контроля над ресурсами и стабильной производительности, такие особенности могут становиться критическими.

    \quad WebAssembly предлагает альтернативную модель выполнения, ориентированную на эффективную доставку и исполнение кода в браузере. В отличие от JavaScript, распространяемого в текстовом виде и требующего многоступенчатой обработки, WebAssembly использует компактное бинарное представление, оптимизированное для быстрой загрузки и декодирования. Модуль WebAssembly компилируется в машинный код практически напрямую, минуя ряд промежуточных этапов, что сокращает время инициализации и снижает вариативность производительности во время выполнения. Такая архитектура делает WebAssembly особенно подходящим для реализации вычислительно интенсивных компонентов веб-приложений, приближая их характеристики к нативным.

\section{Архитектура и модель исполнения}

    \quad Архитектура WebAssembly базируется на концепции виртуальной стековой машины, что обеспечивает компактность инструкций и простоту их валидации. В отличие от традиционных регистровых архитектур, инструкции Wasm оперируют значениями, находящимися в стеке операндов, что позволяет эффективно транслировать код в машинные команды большинства современных CPU. Безопасность исполнения гарантируется строгой статической типизацией: корректность модуля проверяется на этапе декодирования, что полностью исключает возникновение ошибок несоответствия типов (type mismatch) во время выполнения.
    
    \quad Система типов WebAssembly ориентирована на максимальную производительность и включает четыре базовых числовых типа: целые числа (i32, i64) и числа с плавающей запятой (f32, f64). Для задач высокопроизводительных вычислений и линейной алгебры ключевое значение имеет поддержка типа f64, соответствующего стандарту IEEE 754. Это обеспечивает детерминированную семантику арифметических операций и гарантирует воспроизводимость численных результатов вне зависимости от браузера или аппаратной платформы хоста.
    
    \quad Одной из фундаментальных особенностей Wasm является модель линейной памяти. Она представляет собой непрерывный, расширяемый массив сырых байтов, полностью изолированный от управляемой памяти JavaScript. Такая модель напрямую соответствует адресному пространству языков C/C++ и Rust, что позволяет эффективно переносить низкоуровневые алгоритмы, активно использующие указатели и прямые манипуляции с памятью. Взаимодействие с внешней средой (JavaScript) осуществляется через буфер ArrayBuffer, предоставляющий быстрый доступ к данным без необходимости их сложной трансформации.
    
    \quad Структурно модуль WebAssembly организован в виде специализированных секций (типы, функции, глобальные переменные, экспорт/импорт). Такая организация поддерживает технологию потоковой компиляции (streaming compilation), позволяя браузеру начинать генерацию машинного кода еще до завершения загрузки бинарного файла по сети. Весь процесс выполнения протекает в изолированной программной среде («песочнице»). Прямой доступ к системным ресурсам или API браузера ограничен: модуль может взаимодействовать с внешним миром исключительно через явно импортированные функции, что обеспечивает высокий уровень безопасности при сохранении околонативной скорости вычислений.

\section{Многопоточность в браузерной среде}

    \quad Браузерная среда исторически основана на однопоточной модели исполнения: основной поток обрабатывает пользовательские события, выполняет JavaScript-код и управляет отрисовкой интерфейса. Блокировка основного потока длительными вычислениями приводит к потере отзывчивости интерфейса. Механизм Web Workers позволяет создавать фоновые потоки исполнения, каждый из которых работает в изолированном контексте и обменивается данными с основным потоком через передачу сообщений.

    \quad Стандартный обмен данными между потоками основан на копировании, что создаёт существенные накладные расходы для больших числовых массивов. Объект SharedArrayBuffer решает эту проблему, предоставляя разделяемую область памяти, доступную из нескольких потоков одновременно. Синхронизация доступа обеспечивается набором атомарных операций через объект Atomics. В контексте WebAssembly эта архитектура позволяет транслировать вызовы стандарта POSIX pthreads непосредственно в Web Workers, размещая линейную память модуля в SharedArrayBuffer для параллельной обработки данных без избыточного копирования.

    \quad Использование разделяемой памяти в современных браузерах ограничено политикой кросс-доменной изоляции (cross-origin isolation). Техническая реализация доступа к объекту SharedArrayBuffer и высокоточным таймерам требует конфигурации среды окружения на уровне HTTP-протокола. Необходимым условием является установка заголовков Cross-Origin-Embedder-Policy (COEP) и Cross-Origin-Opener-Policy (COOP). Данная настройка формирует изолированный контекст безопасности, позволяющий выполнять многопоточные модули WebAssembly без риска несанкционированного доступа к данным из других контекстов.


\section{Методы компиляции нативного кода в\\ WebAssembly}

    \quad WebAssembly как целевая платформа компиляции поддерживается широким набором инструментариев, различающихся по исходным языкам, архитектуре компилятора и характеристикам генерируемого кода. Выбор конкретного стека определяет архитектуру бинарного модуля, объём сопутствующего рантайма и методику взаимодействия с хост-средой. Рассмотрим основные методы компиляции.

    \quad \textbf{AOT-компиляция (ahead-of-time).} Данный метод предполагает предварительную трансляцию исходного кода непосредственно в бинарный модуль WebAssembly ещё до запуска программы. Именно AOT-компиляция стала основным способом доставки производительных Wasm-модулей в браузер и серверные рантаймы, поскольку она обеспечивает предсказуемое время инициализации, компактный формат поставки и высокий уровень оптимизации. К инструментам этого класса относятся Emscripten, LLVM/Clang, WASI SDK, Rust с \texttt{wasm-bindgen} и \texttt{wasm-pack}, TinyGo, AssemblyScript и другие компиляторные цепочки, генерирующие Wasm напрямую.

\begin{table}[htbp]
\centering
\scriptsize
\renewcommand{\arraystretch}{1.15}
\begin{tabular}{|p{5.8cm}|p{5.8cm}|}
\hline
\textbf{Плюсы} & \textbf{Минусы} \\
\hline
Высокая производительность исполнения & Требуется отдельный этап сборки \\
\hline
Глубокие оптимизации на этапе компиляции & Плохо подходит для динамической генерации кода \\
\hline
Хороший контроль над экспортируемым интерфейсом & Может требовать сложной настройки памяти и рантайма \\
\hline
Удобная интеграция в стандартные сборочные пайплайны & Работа с внешними зависимостями бывает нетривиальной \\
\hline
\end{tabular}
\caption{Преимущества и недостатки AOT-компиляции}
\label{tab:aot_pros_cons}
\end{table}

    \quad \textbf{JIT-компиляция (just-in-time).} При данном методе код компилируется во время исполнения, когда уже известен контекст запуска, профиль нагрузки и характеристики среды. В классическом виде JIT чаще применяется не как основной способ получения Wasm-модуля из исходного текста, а как механизм динамической компиляции промежуточного представления или байткода внутри среды выполнения. В экосистеме WebAssembly этот путь встречается заметно реже, чем AOT, однако важен для рантаймов, где требуется адаптация к данным профилирования и возможность поздней оптимизации. Характерные инструменты и движки этого класса --- Wasmer JIT, Cranelift в JIT-сценариях, а также внутренние JIT-пайплайны движков V8 и SpiderMonkey.

\begin{table}[htbp]
\centering
\scriptsize
\renewcommand{\arraystretch}{1.15}
\begin{tabular}{|p{5.8cm}|p{5.8cm}|}
\hline
\textbf{Плюсы} & \textbf{Минусы} \\
\hline
Оптимизация с учётом реального профиля выполнения & Высокая сложность реализации \\
\hline
Возможность улучшать горячие участки после запуска & Увеличивает время холодного старта \\
\hline
Гибкая адаптация к среде исполнения & Требует более сложного рантайма \\
\hline
Подходит для сценариев поздней оптимизации & Редко используется как основной метод доставки Wasm \\
\hline
\end{tabular}
\caption{Преимущества и недостатки JIT-компиляции}
\label{tab:jit_pros_cons}
\end{table}

    \quad \textbf{Binary translation.} Этот метод основан на переводе уже существующего бинарного кода, байткода или инструкций другой архитектуры в представление, пригодное для исполнения в среде WebAssembly. Такой подход применяется тогда, когда исходный код недоступен либо перенос системы с уровня исходников экономически нецелесообразен. В практических и исследовательских решениях сюда можно отнести CheerpX, некоторые QEMU/LLVM-ориентированные схемы динамической трансляции и другие специализированные системы двоичного переноса.

\begin{table}[htbp]
\centering
\scriptsize
\renewcommand{\arraystretch}{1.15}
\begin{tabular}{|p{5.8cm}|p{5.8cm}|}
\hline
\textbf{Плюсы} & \textbf{Минусы} \\
\hline
Позволяет использовать уже существующий бинарный код & Очень высокая сложность реализации \\
\hline
Не всегда требует доступа к исходникам & Производительность обычно ниже AOT-компиляции \\
\hline
Полезен для переноса legacy-систем & Результат зависит от качества сопоставления архитектур \\
\hline
Может сокращать затраты на переписывание системы & Используется в основном в нишевых сценариях \\
\hline
\end{tabular}
\caption{Преимущества и недостатки binary translation}
\label{tab:binary_translation_pros_cons}
\end{table}

    \quad \textbf{Интерпретация.} В интерпретационном подходе в WebAssembly переносится не сама программа в виде нативно скомпилированного модуля, а интерпретатор, который затем исполняет внешний бинарный код, байткод или сценарный язык инструкция за инструкцией. Этот метод особенно удобен для обеспечения совместимости и быстрого запуска существующих экосистем без радикальной переработки их внутренней модели исполнения. К подобным решениям относятся Pyodide с интерпретатором CPython, Lua- и QuickJS-рантаймы, а также различные эмуляторы и интерпретаторы, скомпилированные в Wasm.

\begin{table}[htbp]
\centering
\scriptsize
\renewcommand{\arraystretch}{1.15}
\begin{tabular}{|p{5.8cm}|p{5.8cm}|}
\hline
\textbf{Плюсы} & \textbf{Минусы} \\
\hline
Высокая гибкость & Производительность ниже AOT-решений \\
\hline
Сравнительно быстрый перенос сложных экосистем & Часто уступает и JIT-подходам \\
\hline
Удобна для задач совместимости & Каждая инструкция проходит дополнительный уровень обработки \\
\hline
Позволяет сохранить полноту поддержки языка & Плохо подходит для тяжёлых вычислительных задач \\
\hline
\end{tabular}
\caption{Преимущества и недостатки интерпретации}
\label{tab:interpretation_pros_cons}
\end{table}

    \quad \textbf{Hybrid-подход.} Гибридная схема объединяет несколько методов компиляции и исполнения в рамках одной системы. На практике это означает, что часть кода может исполняться через интерпретацию, часть --- быть заранее скомпилированной AOT, а наиболее критические участки --- дополнительно оптимизироваться специальным рантаймом. Подобные решения характерны прежде всего для управляемых платформ и больших фреймворков, где необходимо одновременно сохранить совместимость, приемлемую скорость запуска и достаточную производительность. Типичными примерами являются Blazor WebAssembly с AOT-режимом, смешанные пайплайны Mono и ряд managed runtime-решений для Java и .NET.

\begin{table}[htbp]
\centering
\scriptsize
\renewcommand{\arraystretch}{1.15}
\begin{tabular}{|p{5.8cm}|p{5.8cm}|}
\hline
\textbf{Плюсы} & \textbf{Минусы} \\
\hline
Гибкий баланс между скоростью запуска и производительностью & Архитектурно более сложен \\
\hline
Позволяет сочетать сильные стороны разных схем исполнения & Требует большего объёма рантайма \\
\hline
Удобен для крупных управляемых платформ & Сложнее в отладке и сопровождении \\
\hline
Помогает сохранять совместимость платформы & Итоговое поведение зависит сразу от нескольких механизмов \\
\hline
\end{tabular}
\caption{Преимущества и недостатки hybrid-подхода}
\label{tab:hybrid_pros_cons}
\end{table}

    \quad \textbf{VM $\rightarrow$ Wasm.} При данном методе в WebAssembly компилируется виртуальная машина или рантайм, а уже внутри него исполняется байткод исходной платформы. В отличие от более общей интерпретации, здесь обычно речь идёт о полноценных управляемых экосистемах с собственным форматом байткода, системой типов, механизмами загрузки классов или сборки мусора. Этот путь особенно распространён для Java- и .NET-платформ, где требуется сохранить существующий стек библиотек и модель исполнения. Наиболее показательные примеры --- Blazor WebAssembly с рантаймом Mono, Java-ориентированные решения вроде CheerpJ и другие системы, в которых виртуальная машина переносится в браузер как Wasm-модуль.

\begin{table}[htbp]
\centering
\scriptsize
\renewcommand{\arraystretch}{1.15}
\begin{tabular}{|p{5.8cm}|p{5.8cm}|}
\hline
\textbf{Плюсы} & \textbf{Минусы} \\
\hline
Высокая совместимость с исходной экосистемой & Увеличивает размер поставляемого модуля \\
\hline
Можно повторно использовать существующий байткод & Даёт более медленный холодный старт \\
\hline
Сохраняет привычную модель исполнения платформы & Создаёт дополнительные накладные расходы рантайма \\
\hline
Облегчает перенос существующих библиотек & Эффективность зависит от качества реализации виртуальной машины \\
\hline
\end{tabular}
\caption{Преимущества и недостатки метода VM $\rightarrow$ Wasm}
\label{tab:vm_to_wasm_pros_cons}
\end{table}

    \quad Сопоставление перечисленных методов показывает, что между ними существует устойчивый компромисс. Чем ближе метод к прямой предварительной компиляции, тем выше, как правило, производительность и тем меньше объём вспомогательного рантайма. Чем сильнее система опирается на интерпретацию, виртуальную машину или гибридную модель, тем проще сохранить совместимость с исходной платформой, но тем выше накладные расходы на запуск и исполнение. Сравнительные характеристики основных методов приведены в таблице~\ref{tab:compilation_methods}.

\begin{table}[htbp]
\centering
\scriptsize
\renewcommand{\arraystretch}{1.2}
\begin{tabular}{|p{2.4cm}|p{2.0cm}|p{3.2cm}|p{2cm}|p{3.2cm}|}
\hline
\textbf{Метод} & \textbf{Вход} & \textbf{Производительность} & \textbf{Сложность} & \textbf{Использование} \\
\hline
AOT & Исходный код & Высокая & Средняя & Основной способ сборки Wasm-модулей \\
\hline
JIT & Код во время выполнения, IR, байткод & Средняя & Высокая & Редко, в основном во внутренних рантаймах \\
\hline
Binary translation & Готовый бинарь, машинный код, байткод & Средняя или низкая & Очень высокая & Нишевые сценарии переноса legacy-систем \\
\hline
Интерпретация & Бинарь, байткод, сценарный код & Низкая & Низкая или средняя & Совместимость и быстрый перенос экосистем \\
\hline
Hybrid & Смешанный вход & Средняя & Высокая & Managed runtime и крупные платформы \\
\hline
VM $\rightarrow$ Wasm & Байткод виртуальной машины & Средняя & Средняя & Java-, .NET- и другие байткодные экосистемы \\
\hline
\end{tabular}
\caption{Сравнение методов компиляции в WebAssembly}
\label{tab:compilation_methods}
\end{table}


\section{Анализ производительности WebAssembly}

    \quad Производительность WebAssembly зависит не только от самого формата, но и от характера программы, браузера, runtime-среды, компилятора и используемых оптимизаций. В работе [7], где сравнивались 41 C-бенчмарк, 9 пар WebAssembly/JavaScript-программ и 3 реальных приложения, показано, что WebAssembly не всегда автоматически превосходит JavaScript. Современные JIT-компиляторы эффективно оптимизируют JavaScript, тогда как прирост от JIT для WebAssembly в Chrome и Firefox оказался ограниченным. Кроме того, из-за модели линейной памяти WebAssembly в ряде случаев потреблял больше памяти, чем JavaScript [7].

    \quad В работе [8] дополнительно показано различие между микробенчмарками и реальными приложениями: на микробенчмарках WebAssembly в среднем был на 9.84\% медленнее JavaScript и потреблял на 5.18\% больше энергии, однако на реальных приложениях оказался на 17.24\% быстрее и снизил энергопотребление на 24.34\%. При сравнении с нативным кодом важным источником накладных расходов являются проверки границ линейной памяти: в работе [9] показано, что эффективные runtime-среды могут удерживать замедление WebAssembly относительно native на уровне 8--20\%, но результат существенно зависит от реализации runtime.

    \quad Таким образом, WebAssembly следует рассматривать не как универсальный способ автоматического ускорения программ, а как эффективный инструмент для переноса вычислительно нагруженного и заранее оптимизированного кода в веб-среду или изолированную runtime-среду. Наибольший выигрыш достигается в случаях, когда исходный код уже хорошо оптимизирован на уровне алгоритмов и работы с памятью, а компиляция выполняется с учётом особенностей WebAssembly. Для задач линейной алгебры это особенно важно, поскольку итоговая производительность определяется не только самим фактом использования WebAssembly, но и блочной организацией вычислений, локальностью обращений к памяти, качеством компиляции и поддержкой низкоуровневых оптимизаций.

\section{Спецификации BLAS и LAPACK}

    \quad BLAS (Basic Linear Algebra Subprograms) представляет собой спецификацию интерфейсов для базовых операций линейной алгебры. Спецификация организована в три уровня:
    \begin{itemize}
    \item \textbf{Уровень 1} определяет операции вектор-вектор со сложностью $O(n)$: скалярное произведение, вычисление норм, масштабирование, линейные комбинации.
    \item \textbf{Уровень 2} определяет операции матрица-вектор со сложностью $O(n^2)$: умножение матрицы на вектор, решение треугольных систем.
    \item \textbf{Уровень 3} определяет операции матрица-матрица со сложностью $O(n^3)$: матричное умножение, ранговые обновления.
    \end{itemize}
    Операции третьего уровня обладают высоким соотношением вычислений к обращениям в память ($O(n)$ операций на элемент данных), что позволяет эффективно использовать кэш-память и достигать производительности, близкой к пиковой.
        
    \quad LAPACK (Linear Algebra PACKage) представляет собой библиотеку подпрограмм для решения задач линейной алгебры высокого уровня: факторизации матриц, решения систем линейных уравнений, вычисления собственных значений и сингулярного разложения. Подпрограммы LAPACK реализованы через вызовы BLAS, что позволяет автоматически наследовать оптимизации конкретной реализации BLAS для целевой платформы. 

    Соглашение об именовании подпрограмм LAPACK отражает тип данных и характер операции: 
    \begin{enumerate}
        \item первая буква указывает на тип данных (например, \texttt{d} для double precision);
        \item последующие буквы кодируют тип матрицы (\texttt{ge} для общей, \texttt{sy} для симметричной, \texttt{po} для положительно определённой);
        \item окончание определяет выполняемую операцию (\texttt{sv} для решения системы, \texttt{trf} для факторизации).
    \end{enumerate}
    
    \quad OpenBLAS является одной из наиболее эффективных и широко используемых реализаций стандартов BLAS и LAPACK с открытым исходным кодом. Библиотека содержит глубоко оптимизированные вычислительные ядра, написанные на языках C, Fortran и ассемблере, что обеспечивает высокую производительность на различных процессорных архитектурах (x86, ARM, RISC-V и других). Благодаря поддержке многопоточного исполнения и возможности адаптации под специфические ограничения целевых сред, OpenBLAS выступает в качестве промышленного стандарта для интеграции методов высокопроизводительной линейной алгебры в современные программные комплексы и кроссплатформенные системы.
    
\chapter{Проектирование архитектуры платформы}

\section{Требования к платформе и функциональные возможности}

    \quad При проектировании интерактивной веб-платформы основной задачей становится построение архитектуры, сочетающей в себе высокую производительность и модульность с наглядной визуализацией алгоритмов. В отличие от обычных вычислительных программ, такая система ориентирована не только на получение результата, но и на демонстрацию логики вычислительного процесса непосредственно в браузере, что определяет её образовательное назначение - как инструмент изучения численных методов.

    \quad Исходя из этого назначения, к платформе предъявляются следующие функциональные требования:
    \begin{itemize}
        \item решение многомерных систем линейных алгебраических уравнений прямыми и итерационными методами;
        \item решение нелинейных систем уравнений итерационными методами;
        \item построение моделей аппроксимации и регрессии по набору наблюдений;
        \item гибкий ввод данных через матричные сетки, текстовые поля для уравнений и динамические таблицы наблюдений;
        \item пошаговая анимированная визуализация хода алгоритма с синхронной цветовой индикацией обрабатываемых элементов, управлением скоростью воспроизведения и выводом диагностических диаграмм.
    \end{itemize}

    \quad К нефункциональным требованиям относятся: выполнение всех вычислений на стороне клиента, без серверного бэкенда; отсутствие блокировки основного потока браузера при обработке матриц прикладной размерности, что обеспечивает отзывчивость интерфейса; единый пользовательский сценарий взаимодействия для всех классов задач; возможность добавления нового численного метода без модификации существующих модулей.

    \quad Перечисленные требования определяют необходимость разделения платформы на несколько согласованных уровней. \textbf{Интерфейсный уровень} отвечает за ввод данных, навигацию и визуализацию результатов. \textbf{Прикладной уровень} управляет сценариями решения задач и координирует выполнение конкретных численных методов. \textbf{Вычислительный уровень} обеспечивает выполнение ресурсоёмких операций линейной алгебры. Подобная декомпозиция минимизирует связанность подсистем, обеспечивая высокую масштабируемость и упрощая интеграцию новых расчетных модулей(как будет показано в 2.3) .

\section{Выбор технологического стека}

    \quad Проектирование платформы интерактивной визуализации численных алгоритмов требует обоснованного выбора технологий, обеспечивающих выполнение ключевых требований: высокопроизводительные вычисления на стороне клиента, модульная архитектура интерфейса, быстрая начальная загрузка и возможность интерактивной визуализации промежуточных состояний алгоритмов. В данном разделе последовательно рассматриваются решения, принятые для каждого уровня технологического стека.

    \quad Основу вычислительного ядра платформы составляет технология\\ WebAssembly (WASM), выбор которой аргументирован результатами анализа в Главе 1. В качестве базового инструментария компиляции выбран Emscripten SDK, что продиктовано возможностью прямой трансляции высокооптимизированных нативных библиотек линейной алгебры, в частности OpenBLAS и LAPACK, в бинарный формат WASM. Данный подход позволяет использовать промышленно апробированные алгоритмы факторизации матриц и итерационных методов без потери их алгоритмической сложности и эффективности, достигая производительности, сопоставимой с нативными приложениями.

    \quad Для реализации пользовательского интерфейса и управления состоянием приложения выбран фреймворк React.js. Использование компонентно-ориентированного подхода и механизма Virtual DOM позволяет эффективно синхронизировать визуальное представление с результатами вычислений в реальном времени. Интеграция высокопроизводительного WASM-модуля в среду выполнения JavaScript осуществляется посредством асинхронного интерфейса загрузки, что предотвращает блокировку основного потока выполнения (Main Thread) и обеспечивает отзывчивость интерфейса при обработке массивов данных значительной размерности.

    \quad Визуализационный слой платформы базируется на использовании\\ Canvas API. Данный выбор обоснован необходимостью отрисовки динамически изменяющихся многомерных структур, где традиционные методы манипуляции DOM-деревом демонстрируют избыточные вычислительные затраты. Взаимодействие JavaScript-слоя с вычислительным ядром на базе WebAssembly реализовано через управление линейной памятью модуля, что минимизирует накладные расходы на передачу данных и критически важно для обеспечения интерактивности визуализации. При этом инфраструктура SharedArrayBuffer используется для обеспечения многопоточности в вычислительных операциях, формируя когерентную среду, способную решать задачи численного анализа в рамках веб-браузера с соблюдением требований к точности и быстродействию.

\section{Архитектура прототипа}

     \quad В предыдущих разделах были определены требования и технологический стек платформы. Настоящий раздел описывает архитектурные решения, принятые при проектировании: организацию взаимодействия между уровнями, механизм подключения алгоритмов, командный протокол вычислительного ядра и модель пошаговой визуализации.

    \quad Архитектура трёх уровней, обозначенных в разделе 2.1, следует паттерну MVP (Model–View–Presenter). Роль Model выполняет вычислительное ядро: оно инкапсулирует чистые вычисления и не зависит ни от интерфейса, ни от логики конкретного метода. Роль View принадлежит компонентам интерфейсного уровня: они отвечают за ввод данных и отображение результатов, но не содержат алгоритмической логики. Роль Presenter берут на себя модули-солверы прикладного уровня: каждый солвер координирует выполнение метода и управляет визуализацией, выступая единственным посредником между Model и View. Методы, ресурсоёмкие операции которых сводятся к стандартным процедурам линейной алгебры (прямые и итерационные методы для линейных СЛАУ, методы аппроксимации), делегируют их вычислительному ядру. Нелинейные методы (Ньютона, Бройдена, гомотопии, продолжения, простой итерации) реализуются на прикладном уровне без обращения к ядру для основного хода вычислений; для тех из них, что внутренне порождают линейные подсистемы, последние решаются собственными средствами прикладного уровня, а эквивалентная команда полного решения СЛАУ формируется лишь для отображения в журнале шагов в учебных целях. В соответствии с принципом MVP, компоненты интерфейса не обращаются к ядру самостоятельно: страница-контейнер получает абстрактный интерфейс вызова ядра из общего инфраструктурного контекста и передаёт его в контекст солвера, тогда как формирование команд и интерпретация результатов выполняются на стороне прикладного уровня.

\begin{figure}[htbp]
\centering
\resizebox{\textwidth}{!}{%
\begin{tikzpicture}[>=Stealth]

% ---- Interface level ----
\fill[intbg, rounded corners=6pt] (0,0) rectangle (18.4,-3.5);
\draw[intborder, dashed, rounded corners=6pt] (0,0) rectangle (18.4,-3.5);
\node[font=\small\bfseries, text=intborder, anchor=north west] at (0.4,-0.2) {ИНТЕРФЕЙСНЫЙ УРОВЕНЬ};

\node[rectangle, draw=intborder, fill=white, rounded corners=4pt,
      inner sep=5pt, font=\footnotesize, align=center] (sp) at (2.5,-2)
      {Страница СЛАУ\\\scriptsize\color{mutedtext}(линейные системы)};
\node[rectangle, draw=intborder, fill=white, rounded corners=4pt,
      inner sep=5pt, font=\footnotesize, align=center] (nsp) at (6.3,-2)
      {Страница нелинейных\\систем};
\node[rectangle, draw=intborder, fill=white, rounded corners=4pt,
      inner sep=5pt, font=\footnotesize, align=center] (rp) at (10.3,-2)
      {Страница регрессии\\\scriptsize\color{mutedtext}(аппроксимация)};
\node[rectangle, draw=intborder, fill=white, rounded corners=4pt,
      inner sep=5pt, font=\footnotesize, align=center] (ui) at (14.5,-2)
      {форма ввода\\визуализация\\журнал шагов\\индикатор готовности ядра\\\scriptsize\color{mutedtext}(дочерние компоненты страниц)};

% ---- Arrow: interface -> application ----
\draw[->, accentblue, thick] (9.2,-3.5) -- (9.2,-4.6)
      node[midway, right, font=\scriptsize, text=mutedtext] {динамическая загрузка};

% ---- Application level ----
\fill[appbg, rounded corners=6pt] (0,-4.6) rectangle (18.4,-7.8);
\draw[appborder, dashed, rounded corners=6pt] (0,-4.6) rectangle (18.4,-7.8);
\node[font=\small\bfseries, text=applabel, anchor=north west] at (0.4,-4.8) {ПРИКЛАДНОЙ УРОВЕНЬ (солверы-плагины)};

\node[rectangle, draw=appborder, fill=white, rounded corners=4pt,
      inner sep=5pt, font=\footnotesize] (ug) at (1.8,-6.5) {Гаусс};
\node[rectangle, draw=appborder, fill=white, rounded corners=4pt,
      inner sep=5pt, font=\footnotesize] (ul) at (4.1,-6.5) {LU};
\node[rectangle, draw=appborder, fill=white, rounded corners=4pt,
      inner sep=5pt, font=\footnotesize] (un) at (6.6,-6.5) {Ньютон};
\node[rectangle, draw=appborder, fill=white, rounded corners=4pt,
      inner sep=5pt, font=\footnotesize] (us) at (9.2,-6.5) {Сплайн};
\node[rectangle, draw=appborder, fill=white, rounded corners=4pt,
      inner sep=5pt, font=\footnotesize] (dots) at (11.2,-6.5) {\ldots};
\node[rectangle, draw=appborder, fill=white, rounded corners=4pt,
      inner sep=5pt, font=\footnotesize, align=center] (iface) at (15,-6.5)
      {\scriptsize единый интерфейс:\\\scriptsize\color{applabel}\{ метаданные, операция решения \}};

% ---- Arrows: application <-> compute ----
\draw[->, accentblue, thick] (7.5,-7.8) -- (7.5,-9)
      node[midway, right, font=\scriptsize, text=mutedtext] {команда BLAS};
\draw[->, accentblue, thick] (11,-9) -- (11,-7.8)
      node[midway, right, font=\scriptsize, text=mutedtext] {текстовый результат};

% ---- Compute level ----
\fill[compbg, rounded corners=6pt] (0,-9) rectangle (18.4,-12.2);
\draw[compborder, dashed, rounded corners=6pt] (0,-9) rectangle (18.4,-12.2);
\node[font=\small\bfseries, text=complabel, anchor=north west] at (0.4,-9.2) {ВЫЧИСЛИТЕЛЬНЫЙ УРОВЕНЬ};

\node[rectangle, draw=compborder, fill=white, rounded corners=4pt,
      inner sep=5pt, font=\footnotesize, align=center] (wctx) at (3.5,-10.8)
      {\textbf{Адаптер ядра}\\\scriptsize\color{mutedtext}инжектируемый интерфейс};

\node[rectangle, draw=compborder, fill=greenhi, rounded corners=4pt,
      inner sep=5pt, font=\footnotesize, align=center, line width=1.2pt] (blas) at (9.8,-10.8)
      {\textbf{BLAS / LAPACK}\\\scriptsize\color{complabel}библиотека линейной алгебры};

\draw[->, compborder, thick] (wctx) -- (blas);

\node[rectangle, draw=compborder, fill=white, rounded corners=4pt,
      inner sep=5pt, font=\footnotesize, align=center] (ww) at (15.5,-10.8)
      {Пул потоков\\\scriptsize\color{mutedtext}параллельное исполнение};

\end{tikzpicture}%
}
\caption{Трёхуровневая архитектура платформы}
\label{fig:prototype_architecture}
\end{figure}

    \quad Каждый численный метод оформлен как самодостаточный модуль фиксированной структуры --- контракт солвера. Контракт включает общие метаданные метода, набор полей, специфичных для класса задачи, и операцию решения, реализующую полный вычислительный сценарий. Операция решения получает контекст выполнения --- структуру, агрегирующую входные данные задачи, инжектируемый интерфейс вызова ядра, императивные интерфейсы визуализации и журнала шагов, разделяемый флаг пропуска анимации и флаг готовности ядра. Состав контракта приведён в таблице~\ref{tab:solver_contract}.

\begin{table}[htbp]
\centering
\scriptsize
\renewcommand{\arraystretch}{1.15}
\begin{tabular}{|p{4.2cm}|p{1.8cm}|p{6.6cm}|}
\hline
\rowcolor{thbg} \textbf{Поле контракта} & \textbf{Обязат.} & \textbf{Назначение} \\
\hline
\multicolumn{3}{|l|}{\textit{Общие поля (все классы задач)}} \\
\hline
Заголовок и описание метода & да & Отображаются на странице задачи \\
\hline
Уникальный идентификатор & да & Различение методов в общих компонентах визуализации \\
\hline
Базовая задержка между шагами журнала & да & Скорость интерактивной анимации регулируется отдельно \\
\hline
Дополнительные параметры метода & нет & $\epsilon$, максимальное число итераций, релаксационный параметр и т.п. \\
\hline
Операция решения & да & Асинхронная процедура, получающая контекст выполнения \\
\hline
\multicolumn{3}{|l|}{\textit{Линейные СЛАУ и нелинейные системы (общие)}} \\
\hline
Размерности задачи и предзаполненного примера & да & \\
\hline
\multicolumn{3}{|l|}{\textit{Линейные СЛАУ}} \\
\hline
Предзаполненный пример & нет & Матрица коэффициентов и вектор правой части \\
\hline
Подпись поля ввода матрицы & нет & Используется при дополнительных требованиях к матрице \\
\hline
\multicolumn{3}{|l|}{\textit{Нелинейные системы}} \\
\hline
Предзаполненная система уравнений & нет & \\
\hline
Начальное приближение & нет & \\
\hline
\multicolumn{3}{|l|}{\textit{Аппроксимация и регрессия}} \\
\hline
Число входных признаков & нет & \\
\hline
Предзаполненный набор наблюдений & нет & \\
\hline
Минимально допустимое число точек & нет & По умолчанию определяется методом \\
\hline
\end{tabular}
\caption{Контракт модуля-солвера}
\label{tab:solver_contract}
\end{table}

    \quad Страницы-контейнеры загружают модуль солвера динамически и передают его конфигурацию в стандартный набор компонентов. Добавление нового метода не требует модификации существующего кода: достаточно создать модуль с заданной структурой и зарегистрировать его в каталоге методов.

    \quad Взаимодействие прикладного уровня с вычислительным ядром организовано как текстовый командный протокол. Солвер формирует строку с именем BLAS- или LAPACK-подпрограммы и её параметрами, передаёт её через инжектируемый интерфейс и получает текстовый результат. Вся логика управления памятью ядра, кодирования команды и перехвата вывода инкапсулирована в адаптере и скрыта от солвера. Классы операций, используемые при реализации методов, приведены в таблице~\ref{tab:blas_commands}.

\begin{table}[htbp]
\centering
\scriptsize
\renewcommand{\arraystretch}{1.15}
\begin{tabular}{|p{6.0cm}|p{6.5cm}|}
\hline
\rowcolor{thbg} \textbf{Класс операции} & \textbf{Назначение в контексте алгоритма} \\
\hline
Поиск индекса экстремального элемента вектора & Выбор ведущего элемента \\
\hline
Линейная комбинация векторов & Элиминация при прямом ходе \\
\hline
Перестановка векторов & Перестановка строк при выборе ведущего элемента \\
\hline
Решение треугольной системы & Обратная подстановка \\
\hline
Полное решение СЛАУ & Независимая верификация результата \\
\hline
Решение симметрично-положительно-определённой СЛАУ & Метод Холецкого \\
\hline
Матрично-векторное произведение & Вычисление невязки и итерационных шагов \\
\hline
Скалярное произведение и норма & Контроль сходимости итерационных методов \\
\hline
\end{tabular}
\caption{Классы операций вычислительного ядра, используемые в методах}
\label{tab:blas_commands}
\end{table}

    \quad В отличие от подхода, при котором вычислительное ядро выполняет метод целиком и возвращает финальный результат, в данной архитектуре алгоритм реализован как последовательность элементарных операций на прикладном уровне. Каждая операция может включать один или несколько вызовов ядра, обновление визуального состояния и запись в журнал. Между операциями вставляются асинхронные паузы, которые передают управление браузеру и позволяют отрисовать промежуточное состояние до продолжения вычислений. Компоненты визуализации и журнала предоставляют солверу императивный интерфейс: методы для подсветки ячеек матрицы, поэлементного обновления значений, добавления записей с описанием шага и соответствующей BLAS-командой. Пользователь может выбирать скорость воспроизведения или пропустить анимацию, при этом алгоритм доводит вычисления до конца и отрисовывает финальное состояние.

    \quad На интерфейсном уровне выделены три типа расчётных страниц-контейнеров, каждый из которых спроектирован под свой класс задач (помимо них в платформе присутствуют также вспомогательные навигационные страницы --- главная и каталог методов). Страница линейных СЛАУ предоставляет матричную сетку $n \times n$ с вектором правой части и визуализирует процесс через интерактивную таблицу с подсветкой. Страница нелинейных систем использует текстовые поля для ввода уравнений и вектор начального приближения, отображая итерации в виде таблицы приближений. Страница регрессии предлагает динамическую таблицу наблюдений и визуализирует результат набором диаграмм. При этом все три типа страниц реализуют единый пользовательский сценарий: ввод и валидация на стороне формы $\rightarrow$ проверка готовности ядра $\rightarrow$ формирование контекста $\rightarrow$ вызов операции решения $\rightarrow$ отображение результата. Алгоритмическая логика полностью инкапсулирована в модуле-солвере.

    \quad Прямые методы решения линейных систем (Гаусса, LU, QR) поддерживают автоматический выбор между двумя режимами вычислений. Если все элементы входной матрицы и вектора правой части являются целыми числами, шаги алгоритма --- выбор ведущего элемента, элиминация, обратная подстановка --- выполняются в поле рациональных дробей на прикладном уровне, сохраняя точность на каждом промежуточном шаге. В противном случае используется арифметика с плавающей запятой, а перечисленные операции делегируются вычислительному ядру. В обоих режимах после получения решения выполняется независимая верификация через полное решение исходной системы средствами ядра. Точный режим особенно ценен в образовательном контексте: промежуточные результаты отображаются в виде обыкновенных дробей, что позволяет проверить каждый шаг вручную.


\section{Структуры данных и интерфейсы взаимодействия}

    \quad Данный раздел дополняет описание архитектуры со стороны организации данных. Основная задача проектирования в рамках этого контура заключается в определении форматов представления многомерных величин на каждом уровне и поиске механизмов их эффективной трансформации. Рассматриваются промежуточные структуры, необходимые для пошаговой визуализации, а также способы обеспечения согласованности данных между вычислительным ядром и графической подсистемой.

    \vspace{0.5em}

    \textbf{\large Логическое и линейное представление матриц}

    \vspace{0.5em}

    \quad Между прикладным уровнем и вычислительным ядром данные пересекают архитектурную границу, на которой меняются их требования к памяти. На прикладном уровне естественной формой матрицы является массив строк, где каждая строка - отдельный массив чисел. Такое представление удобно для ручного ввода, поэлементного обновления, отрисовки в виде таблицы и отслеживания изменений конкретной ячейки в анимации. Спецификации BLAS и LAPACK, наоборот, предполагают непрерывное размещение элементов в одном блоке памяти, поскольку именно непрерывность обеспечивает эффективное использование кэш-памяти и блочных оптимизаций.

    \quad В архитектуре проектируется два согласованных представления матрицы:

    \begin{itemize}
        \item \emph{Логическое (двумерное)} : существует на прикладном уровне и в журнале шагов. Используется для ввода, валидации, пошаговых преобразований и подсветки ячеек.
        \item \emph{Линейное (одномерное)} : построчная (row-major) развёртка матрицы $n \times n$ в  последовательность длины $n^2$. Формируется только в момент вызова вычислительного ядра и используется как часть строки команды, передаваемой в линейную память WebAssembly.
    \end{itemize}

    \quad Векторы во всех контекстах остаются одномерными и не требуют преобразования при передаче в ядро. Преобразование двумерного представления в линейное выполняется на стороне солвера непосредственно перед формированием команды, что позволяет хранить матрицы в удобной форме на всём остальном протяжении расчёта.

    \vspace{0.5em}

    \textbf{\large Типы рабочих данных}

    \vspace{0.5em}

    \quad Помимо самих матриц и векторов, в системе проектируется ряд вспомогательных типов данных, специфичных для разных классов задач. Сводка этих типов приведена в таблице~\ref{tab:working_types}.

\begin{table}[htbp]
\centering
\scriptsize
\renewcommand{\arraystretch}{1.15}
\begin{tabular}{|p{4.0cm}|p{8.5cm}|}
\hline
\rowcolor{thbg} \textbf{Тип данных} & \textbf{Назначение} \\
\hline
Рациональное число & Пара числитель-знаменатель для режима точной арифметики; служит элементом матрицы и вектора при целочисленных входных данных. \\
\hline
Снимок шага & Запись текущего состояния алгоритма, формируемая между элементарными операциями. \\
\hline
Итерационная запись & Кортеж номера итерации, текущего приближения и значения нормы для отображения в журнале итерационных методов. \\
\hline
Набор наблюдений & Список строк \emph{(признаки, значения)} для задач регрессии и аппроксимации. \\
\hline
Спецификация подсветки & Метаструктура с координатами выделяемых строк, столбцов и клеток, сопровождающая каждый снимок матрицы. \\
\hline
\end{tabular}
\caption{Вспомогательные типы рабочих данных}
\label{tab:working_types}
\end{table}

    \quad Поскольку все числовые операции солвера выполняются через единый набор функций, переключение между рациональным и вещественным режимом происходит на уровне выбора реализации этих операций, а сам алгоритм работает с одной и той же структурой данных. Это устраняет дублирование логики между двумя режимами и позволяет точному режиму использовать тот же набор пошаговых снимков, что и режим с плавающей запятой.

    \vspace{0.5em}

    \textbf{\large Состав снимка шага}

    \vspace{0.5em}

    \quad Реализация алгоритмов в виде последовательности атомарных операций (см. разд. 2.3) предполагает фиксацию промежуточных состояний системы. Таким состоянием является снимок шага — самодостаточный объект данных, передаваемый от вычислительного ядра к слою визуализации. В его состав входят:

    \begin{itemize}
        \item заголовок и описание операции в терминах линейной алгебры (например, элементарные преобразования строк или итерационное уточнение);
        \item текущее состояние объектов (матриц и векторов) в логическом представлении;
        \item метаданные подсветки: индексы ведущих, модифицируемых или целевых элементов (например, треугольных факторов);
        \item параметры сходимости (для итерационных методов): номер итерации, текущее приближение и норма невязки;
        \item строковая интерпретация команды ядра для отображения в техническом журнале.
    \end{itemize}

    \quad Снимок проектируется как независимая структура, содержащая полный контекст для отрисовки. Это исключает обращение компонентов визуализации к глобальному состоянию алгоритма и позволяет синхронно обновлять различные представления — от анимированных таблиц на Canvas до текстового лога выполнения.

    \vspace{0.5em}

    \textbf{\large Структуры результатов по классам задач}

    \vspace{0.5em}

    \quad Для обеспечения универсальности платформы структуры итоговых данных дифференцированы по типам решаемых математических задач. Каждая структура спроектирована таким образом, чтобы содержать исчерпывающую информацию для визуализации и верификации без необходимости повторного обращения к вычислительному ядру.

    \quad В задачах СЛАУ результат представляет собой агрегат, включающий вектор решения, факторы разложения (в случае использования прямых методов) и блок верификации. Наличие в структуре независимо вычисленного эталонного решения и оценки невязки обеспечивает встроенный контроль корректности расчётов и позволяет пользователю мгновенно оценить точность полученного результата.

    \quad При работе с нелинейными системами результатом является массив итерационных состояний. Структура хранит последовательность пар «приближение - норма невязки», флаг сходимости, итоговое число выполненных итераций и детерминированную причину останова, такую как достижение заданной точности или исчерпание лимита итераций. Данная организация данных оптимизирована для автоматического формирования графиков сходимости и итерационных таблиц.

    \quad Для задач регрессии и аппроксимации используется композитная структура, разделенная на два функциональных блока. Аналитический блок содержит коэффициенты модели и статистические метрики качества, а графический блок включает предвычисленные значения функции на наборе узлов и данные для построения диаграмм рассеяния. Такое разделение позволяет одному объекту результата обслуживать несколько различных типов визуализаций, сохраняя при этом полный аналитический контекст расчёта.

    \vspace{0.5em}

    \textbf{\large Интерфейсы взаимодействия между уровнями}

    \vspace{0.5em}

    \quad Спроектированные структуры данных функционируют в рамках четырех интерфейсов, фиксирующих границы между уровнями системы (табл. \ref{tab:interfaces}).

\begin{table}[htbp]
\centering
\scriptsize
\renewcommand{\arraystretch}{1.15}
\begin{tabular}{|p{3.4cm}|p{2.6cm}|p{3.0cm}|p{3.5cm}|}
\hline
\rowcolor{thbg} \textbf{Интерфейс} & \textbf{Направление} & \textbf{Форма обмена} & \textbf{Режим} \\
\hline
Страница~--~солвер & страница~$\rightarrow$~солвер & объект контекста вызова & асинхронный, однократный \\
\hline
Контракт солвера & солвер~$\rightarrow$~страница & метаданные и \texttt{solve(ctx)} & декларативный \\
\hline
Командный протокол ядра & солвер~$\leftrightarrow$~ядро & текстовая команда / текстовый ответ & синхронный, без состояния \\
\hline
Алгоритм~--~визуализация & солвер~$\rightarrow$~слой отрисовки & журнал снимков и анимационные примитивы & потоковый, с прерыванием \\
\hline
\end{tabular}
\caption{Интерфейсы взаимодействия между уровнями архитектуры}
\label{tab:interfaces}
\end{table}

    \quad \textbf{Интерфейс <<страница-солвер>>} инкапсулирует параметры задачи в единый объект контекста. Это изолирует логику ввода от математических алгоритмов: страница управляет состоянием формы и дескрипторами Canvas, не участвуя в формировании команд ядра.

    \quad \textbf{Контракт модуля-солвера} унифицирует интеграцию методов через единый набор метаданных и исполняемую функцию \texttt{solve}. Декларативная часть контракта позволяет странице автоматически генерировать формы ввода, что упрощает масштабирование библиотеки без изменения ядра системы.

    \quad \textbf{Командный протокол вычислительного ядра} формализует вызовы WebAssembly-модуля. Солвер транслирует операции в текстовые команды (стандарты CBLAS/LAPACKE), включая конфигурационные флаги и линейные массивы. Текстовый формат обеспечивает прозрачность логирования и упрощает отладку взаимодействия с низкоуровневой памятью.

    \quad \textbf{Интерфейс <<алгоритм-визуализация>>} связывает расчет с графическим слоем через два канала: журнал снимков и анимационные примитивы. Журнал гарантирует воспроизводимость состояний для экспорта, а примитивы обеспечивают динамическую подсветку элементов в процессе вычисления.

    \quad Предложенная схема интерфейсов делит платформу на слабо связанные слои. Это позволяет независимо модифицировать способы взаимодействия с пользователем, заменять вычислительные сборки ядра или внедрять новые типы графических отчетов без нарушения целостности системы.






\chapter{Реализация прототипа}


    \quad Настоящая глава посвящена программной реализации архитектурных решений, описанных в Главе 2. Основное внимание уделено сборке и подключению вычислительного ядра к интерфейсу, программной структуре солверов и реализации алгоритмов точной арифметики. Также приводится описание графической подсистемы и принципов взаимодействия между страницами-контейнерами и слоем визуализации.. 

\section{Подготовка вычислительного ядра}

    \quad Ядро реализовано как командная оболочка над OpenBLAS и LAPACK. Основной модуль  \texttt{shell\_cblas.c}  определяет таблицу диспетчеризации \texttt{g\_commands} и единную точку входа \texttt{run\_command}, предназначенную для обработки текстовых команд. В соответствии с принятым протоколом (см. разд. 2.3), диспетчер выполняет токенизацию входной строки, идентификацию команды и передачу управления профильному обработчику \texttt{cmd\_*}. В системе зарегистрировано 44 команды, охватывающие три уровня BLAS, ключевые подпрограммы LAPACK и операции для разреженных матриц (CSR). Распределение функций по категориям представлено в таблице~\ref{tab:core_commands}. 

\begin{table}[htbp]
\centering
\scriptsize
\renewcommand{\arraystretch}{1.15}
\begin{tabular}{|p{2.3cm}|c|p{8.7cm}|}
\hline
\rowcolor{thbg} \textbf{Категория} & \textbf{N} & \textbf{Команды} \\
\hline
BLAS Level~1 & 10 & \texttt{ddot}, \texttt{dnrm2}, \texttt{dasum}, \texttt{idamax}, \texttt{dscal}, \texttt{daxpy}, \texttt{dcopy}, \texttt{dswap}, \texttt{drotg}, \texttt{drot} \\
\hline
BLAS Level~2 & 7 & \texttt{dgemv}, \texttt{dtrsv}, \texttt{dger}, \texttt{dsymv}, \texttt{dsyr}, \texttt{dsyr2}, \texttt{dtrmv} \\
\hline
BLAS Level~3 & 5 & \texttt{dgemm}, \texttt{dtrsm}, \texttt{dtrmm}, \texttt{dsyrk}, \texttt{dsymm} \\
\hline
LAPACK & 17 & \texttt{dgesv}, \texttt{dgetrf}, \texttt{dgetrs}, \texttt{dpotrf}, \texttt{dpotrs}, \texttt{dposv}, \texttt{dgeqrf}, \texttt{dorgqr}, \texttt{dormqr}, \texttt{dgels}, \texttt{dgelsd}, \texttt{dgelss}, \texttt{dgtsv}, \texttt{dptsv}, \texttt{dgbsv}, \texttt{dsyev}, \texttt{dgesvd} \\
\hline
Sparse (CSR) & 3 & \texttt{csr\_spmv}, \texttt{csr\_spmm}, \texttt{csr\_to\_dense} \\
\hline
Служебные & 2 & \texttt{threads}, \texttt{test} \\
\hline
\end{tabular}
\caption{Категории команд вычислительного ядра}
\label{tab:core_commands}
\end{table}

    \quad Для обеспечения отказоустойчивости ядра реализован механизм перехвата исключительных ситуаций. Парсинг параметров выполняется функциями \texttt{toint} и \texttt{todbl} с валидацией типов и диапазона. В случае возникновения критической ошибки внутри обработчика (некорректный аргумент, нехватка памяти, аномальный размер) управление передается через \texttt{longjmp} в точку возврата \texttt{g\_cmd\_abort}, установленную диспетчером. Такой подход предотвращает аварийное завершение WASM-модуля и позволяет вернуть в JavaScript-слой диагностическое сообщение. Результаты вычислений транслируются в стандартный вывод через \texttt{printf}, который перенаправляется окружением Emscripten в асинхронный поток данных приложения.

    \quad Процесс сборки автоматизирован скриптом \texttt{build.sh}. Библиотека\\ OpenBLAS компилируется под целевую платформу в режиме \texttt{TARGET=GENERIC} с поддержкой многопоточности (\texttt{USE\_THREAD=1}, \texttt{NUM\_THREADS=4}).Компоновка модуля  \texttt{shell\_cblas.c} с объектным файлом \texttt{libopenblas.a} выполняется компилятором \texttt{emcc} с оптимизацией \texttt{-O2}. Ключевые параметры компоновщика, определяющие конфигурацию среды исполнения, приведены в таблице~\ref{tab:emcc_flags}. Итоговые артефакты включают бинарный модуль \texttt{shell\_cblas.wasm}, связующий JavaScript-код и вспомогательный поток \texttt{shell\_cblas.worker.js} для реализации pthread.

\begin{table}[htbp]
\centering
\scriptsize
\renewcommand{\arraystretch}{1.15}
\begin{tabular}{|p{5.6cm}|p{7.4cm}|}
\hline
\rowcolor{thbg} \textbf{Параметр} & \textbf{Назначение} \\
\hline
\texttt{-pthread}, \texttt{PTHREAD\_POOL\_SIZE=4} & Предсозданный пул из четырёх Web Workers под pthreads OpenBLAS \\
\hline
\texttt{INITIAL\_MEMORY=67108864} & Начальный размер линейной памяти WASM (64~МБ) \\
\hline
\texttt{ALLOW\_MEMORY\_GROWTH=1} & Динамическое расширение памяти при необходимости \\
\hline
\texttt{INVOKE\_RUN=0}, \texttt{EXIT\_RUNTIME=0} & Ручной запуск \texttt{main}, рантайм не завершается после возврата \\
\hline
\texttt{EXPORTED\_FUNCTIONS} & \texttt{\_main}, \texttt{\_run\_command}, \texttt{\_malloc}, \texttt{\_free} \\
\hline
\texttt{EXPORTED\_RUNTIME\_METHODS} & \texttt{stringToUTF8}, \texttt{lengthBytesUTF8} (доступ из JS-обёртки) \\
\hline
\texttt{SUPPORT\_LONGJMP=emscripten} & Поддержка \texttt{longjmp} для механизма обработки ошибок \\
\hline
\end{tabular}
\caption{Ключевые параметры компоновщика \texttt{emcc}}
\label{tab:emcc_flags}
\end{table}
\section{Подключение ядра к интерфейсному уровню}

     \quad Интеграция вычислительного модуля реализована в провайдере \texttt{WasmProvider} (\texttt{WasmContext.jsx}), который при монтировании инициализирует глобальный объект \texttt{window.Module}, задающий параметры рантайма Emscripten. В конфигурации задействованы три ключевых поля: функция \texttt{locateFile} обеспечивает корректную маршрутизацию запросов к бинарному файлу \texttt{shell\_cblas.wasm} в целевой каталог проекта, поле \texttt{print} реализует прокси-функцию для перехвата и перенаправления стандартного вывода ядра, а обработчик\\ \texttt{onRuntimeInitialized} сигнализирует о полной готовности среды исполнения путем выставления реактивного флага \texttt{wasmReady}. Такой подход позволяет централизованно управлять жизненным циклом WebAssembly-модуля и связывать низкоуровневый вывод C-кода с состоянием React-приложения.

  \quad Особенностью рантайма Emscripten является фиксация ссылки на метод \texttt{Module.print} в момент загрузки. Для обеспечения динамической смены потребителя данных реализован уровень косвенности: в замыкании контекста объявлена переменная-указатель \texttt{\_printHandler}, а стабильный метод \texttt{Module.print} лишь транслирует аргументы в текущую активную функцию. Это позволяет перехватывать вывод ядра в любой момент жизненного цикла приложения простым переназначением \texttt{\_printHandler}.

    \quad На данном механизме базируется функция \texttt{runBlas(cmd)}, реализующая текстовый протокол взаимодействия (см. разд. 2.3). С точки зрения прикладного уровня вызов \texttt{runBlas} является синхронным и атомарным, а вся низкоуровневая работа с памятью инкапсулирована внутри функции. Алгоритм её работы представлен в таблице~\ref{tab:runblas_pipeline}.
    
\begin{table}[htbp]
\centering
\scriptsize
\renewcommand{\arraystretch}{1.15}
\begin{tabular}{|c|p{11cm}|}
\hline
\rowcolor{thbg} \textbf{Шаг} & \textbf{Действие} \\
\hline
1 & Сохранить текущий \texttt{\_printHandler} и заменить его на накопитель строк \\
\hline
2 & Вычислить длину UTF-8-представления команды функцией \texttt{lengthBytesUTF8} \\
\hline
3 & Выделить буфер в линейной памяти WASM через \texttt{Module.\_malloc} \\
\hline
4 & Записать строку в буфер функцией \texttt{stringToUTF8} \\
\hline
5 & Вызвать экспортированную точку \texttt{Module.\_run\_command} с указателем на буфер \\
\hline
6 & Освободить буфер через \texttt{Module.\_free}, восстановить прежний \texttt{\_printHandler}, вернуть накопленный текст \\
\hline
\end{tabular}
\caption{Последовательность операций функции \texttt{runBlas}}
\label{tab:runblas_pipeline}
\end{table}

    \quad Контекст экспортирует специализированныый хук \texttt{useWasm}, предоставляющий компонентам реактивный доступ к состоянию \texttt{wasmReady} и функции \texttt{runBlas}. Это позволяет интерфейсным элементам вызывать методы линейной алгебры так же, как обычные функции JavaScript, полностью игнорируя специфику работы с WebAssembly.


\section{Модули-солверы}

    \quad Численные методы размещены в каталоге \texttt{src/hooks/solv\-ers/}, разделённом на подкаталоги \texttt{linear/}, \texttt{nonlinear/} и \texttt{approx/} по классу решаемой задачи. Каждый модуль экспортирует объект, удовлетворяющий контракту солвера (табл.~\ref{tab:solver_contract}). Номенклатура реализованных модулей по классам приведена в таблицах~\ref{tab:solver_inventory_linear}--\ref{tab:solver_inventory_approx}. Значения столбца <<Ядро>>: \textit{да}~--- ядро используется в основном ходе вычислений; \textit{верификация}~--- ядро применяется только для контрольной проверки решения; \textit{журнал}~--- ядро вызывается лишь для формирования эквивалентной BLAS-команды и независимой проверки.

\begin{table}[H]
\centering
\scriptsize
\renewcommand{\arraystretch}{1.15}
\begin{tabular}{|p{3.5cm}|p{3.0cm}|c|c|}
\hline
\rowcolor{thbg} \textbf{Метод} & \textbf{Файл} & \textbf{Ядро} & \textbf{Точн. реж.} \\
\hline
Метод Гаусса & \texttt{useGauss.js} & да & да \\ \hline
LU-разложение & \texttt{useLU.js} & да & да \\ \hline
QR-разложение & \texttt{useQR.js} & верификация & да \\ \hline
Метод Холецкого & \texttt{useCholesky.js} & да & нет \\ \hline
Метод Якоби & \texttt{useJacobi.js} & да & нет \\ \hline
Метод Зейделя & \texttt{useSeidel.js} & да & нет \\ \hline
SOR & \texttt{useSOR.js} & да & нет \\ \hline
MINRES & \texttt{useMinRes.js} & да & нет \\ \hline
BiCG & \texttt{useBiCG.js} & да & нет \\ \hline
GMRES & \texttt{useGMRES.js} & да & нет \\ \hline
\end{tabular}
\caption{Модули-солверы для линейных СЛАУ (\texttt{linear/})}
\label{tab:solver_inventory_linear}
\end{table}

\begin{table}[H]
\centering
\scriptsize
\renewcommand{\arraystretch}{1.15}
\begin{tabular}{|p{3.5cm}|p{3.0cm}|c|c|}
\hline
\rowcolor{thbg} \textbf{Метод} & \textbf{Файл} & \textbf{Ядро} & \textbf{Точн. реж.} \\
\hline
Метод Ньютона & \texttt{useNewton.js} & журнал & нет \\ \hline
Метод Бройдена & \texttt{useBroyden.js} & журнал & нет \\ \hline
Простая итерация & \texttt{useIteration.js} & нет & нет \\ \hline
Гомотопия & \texttt{useHomotopy.js} & нет & нет \\ \hline
Продолжение по параметру & \texttt{useContinuation.js} & нет & нет \\ \hline
\end{tabular}
\caption{Модули-солверы для нелинейных систем (\texttt{nonlinear/})}
\label{tab:solver_inventory_nonlinear}
\end{table}

\begin{table}[H]
\centering
\scriptsize
\renewcommand{\arraystretch}{1.15}
\begin{tabular}{|p{3.5cm}|p{3.0cm}|c|c|}
\hline
\rowcolor{thbg} \textbf{Метод} & \textbf{Файл} & \textbf{Ядро} & \textbf{Точн. реж.} \\
\hline
Линейная регрессия & \texttt{useLinearRegression.js} & да & нет \\ \hline
Полиномиальная регрессия & \texttt{usePolyRegression.js} & да & нет \\ \hline
Метод наим.\ квадратов & \texttt{useLeastSquares.js} & да & нет \\ \hline
RBF-интерполяция & \texttt{useRBF.js} & да & нет \\ \hline
Кубический сплайн & \texttt{useSpline.js} & да & нет \\ \hline
\end{tabular}
\caption{Модули-солверы для аппроксимации и регрессии (\texttt{approx/})}
\label{tab:solver_inventory_approx}
\end{table}

    \quad Маршруты регистрируются в \texttt{src/App.jsx}: словари \texttt{solverConfigs}, \texttt{nonlinear\-Configs} и \texttt{approxConfigs} сопоставляют ключу маршрута динамический импорт соответствующего модуля. Страницы загружаются через\\ \texttt{React.lazy} внутри \texttt{<Suspense>}, что обеспечивает разбиение бандла: модуль каждого солвера выделяется в отдельный chunk и подгружается только при переходе на маршрут.

    \quad Двухфазная схема (см.\ разд.~2.3) реализуется служебными точками. В интерактивной фазе между шагами вставляется ожидание\\ \texttt{animSleep(viz.getSpeed(), skipRef)}; нажатие кнопки <<Пропустить>> устанавливает \texttt{skipRef.current = true}, после чего ветви алгоритма сворачиваются до окончательного состояния. В журнальной фазе те же шаги исполняются без задержек, а каждая операция фиксируется парой \texttt{stepLog.addStep(\ldots)} и \texttt{await stepLog.showStep(\ldots)}.

    \quad Композиция операций ядра по классам методов приведена в таблице~\ref{tab:method_blas}.

\begin{table}[!htbp]
\centering
\scriptsize
\renewcommand{\arraystretch}{1.15}
\begin{tabular}{|p{4.2cm}|p{8.5cm}|}
\hline
\rowcolor{thbg} \textbf{Класс методов} & \textbf{Используемые операции ядра} \\
\hline
Гаусс, LU & \texttt{idamax}, \texttt{dswap}, \texttt{daxpy}, \texttt{dtrsv}; верификация через \texttt{dgesv} \\
\hline
QR (Хаусхолдер) & только верификация решения через \texttt{dgesv} \\
\hline
Холецкий & \texttt{dpotrf}, \texttt{dpotrs}; верификация через \texttt{dposv} \\
\hline
Итерационные методы (Якоби, Зейдель, SOR, BiCG, MINRES, GMRES) & \texttt{dgemv}, \texttt{daxpy}, \texttt{ddot}; верификация через \texttt{dgesv} \\
\hline
Ньютон, Бройден & внутреннее решение линеаризации через \texttt{luSolve} (JS); в журнал помещается эквивалентная команда \texttt{dgesv} для независимой проверки \\
\hline
Регрессия и аппроксимация & решение нормальных уравнений через \texttt{dgesv} \\
\hline
\end{tabular}
\caption{Композиция операций ядра по классам методов}
\label{tab:method_blas}
\end{table}

    \quad Прямые методы Гаусса и LU реализуют поэтапную элиминацию вызовами \texttt{idamax} (выбор ведущего элемента), \texttt{dswap} (перестановка строк), \texttt{daxpy} (обнуление поддиагональных элементов) и \texttt{dtrsv} (обратная подстановка). QR-разложение выполняется в JS методом отражений Хаусхолдера и обращается к ядру лишь для верификации; такой выбор обусловлен отсутствием в развёрнутой сборке OpenBLAS точек \texttt{dgeqrf}/\texttt{dorgqr}. Метод Холецкого опирается на пару LAPACK-вызовов \texttt{dpotrf}/\texttt{dpotrs} с проверкой через \texttt{dposv}. Итерационные методы выражают каждый внутренний шаг через \texttt{dgemv}, \texttt{daxpy} и \texttt{ddot}. Методы Ньютона и Бройдена решают линеаризацию внутренней процедурой \texttt{luSolve}, а в журнал шагов помещается эквивалентная команда \texttt{dgesv}, результат которой сравнивается с собственным для контроля. Регрессионные методы сводят задачу к нормальным уравнениям и решают их единственным вызовом \texttt{dgesv}. В \texttt{useGauss.js} дополнительно зарегистрирован обработчик\\ \texttt{viz.enableBacksub(\ldots)}, реализующий развёртку обратной подстановки в стиле Гаусса--Жордана; пользователь активирует её отдельной кнопкой.

    \quad Метод-специфичные параметры (точность $\varepsilon$, максимальное число итераций, релаксационный коэффициент $\omega$ в SOR, шаг $\tau$ в гомотопии и т.\,п.) передаются через поле \texttt{extra\-Params} контракта; их значения собираются компонентом \texttt{MatrixInput} и попадают в контекст солвера под именем \texttt{data.extra}.

\section{Точная арифметика рациональных дробей}

    \quad Точный режим реализован классом \texttt{Frac} в \texttt{src/utils/solverUtils.js}: пара целочисленных полей \texttt{n} (числитель) и \texttt{d} (знаменатель). Конструктор приводит знаменатель к положительному значению и сокращает дробь по наибольшему общему делителю, вычисляемому алгоритмом Евклида. Полный набор операций класса приведён в таблице~\ref{tab:frac_ops}.

\begin{table}[!htbp]
\centering
\scriptsize
\renewcommand{\arraystretch}{1.15}
\begin{tabular}{|p{4.2cm}|p{3.5cm}|p{5.0cm}|}
\hline
\rowcolor{thbg} \textbf{Метод} & \textbf{Операция} & \textbf{Возвращаемый результат} \\
\hline
\texttt{add(o)}, \texttt{sub(o)}, \texttt{mul(o)}, \texttt{div(o)} & арифметика & новый объект \texttt{Frac} \\ \hline
\texttt{neg()}, \texttt{abs()} & унарные & новый объект \texttt{Frac} \\ \hline
\texttt{isZero()} & предикат & булево значение \\ \hline
\texttt{absVal()} & $|n/d|$ & число с плавающей запятой \\ \hline
\texttt{toNumber()} & $n/d$ & число с плавающей запятой \\ \hline
\texttt{toString()} & сериализация & строка \texttt{n/d}, либо \texttt{n} при $d=1$ \\ \hline
\end{tabular}
\caption{Операции класса \texttt{Frac}}
\label{tab:frac_ops}
\end{table}

    \quad Переключение режимов опирается на предикат \texttt{isAllInt(A, b)}: при целочисленных $A$ и $b$ входные данные конвертируются функциями \texttt{toFracMatrix} и \texttt{toFracVec}. Сам алгоритм солвера сохраняется единой функцией; различаются лишь определения арифметических операций, выбираемые тернарной развязкой по флагу \texttt{exact} (в типичной форме \texttt{const add = exact ? (a,b)\\ => a.add(b) : (a,b) => a + b}). Такой приём устраняет дублирование между режимами и позволяет точному режиму использовать ту же последовательность шагов визуализации, что и численному.

    \quad Для типографского отображения дробей в HTML предусмотрена функция \texttt{fmtNumHtml}: при $d \ne 1$ она формирует разметку\\ \texttt{<span class="frac"><sup>n</sup>\&frasl;<sub>d</sub></span>} с учётом знака. Поскольку ядро работает только с двойной точностью, верификация в режиме \texttt{Frac} выполняется по числовым значениям, полученным методом \texttt{toNumber}. Точный режим поддержан в трёх прямых методах решения СЛАУ: Гаусса, LU и QR.

\section{Слой визуализации}

    \quad Компоненты слоя визуализации построены как React-компоненты с императивным интерфейсом через \texttt{ref}. Такая схема выбрана для того, чтобы солвер мог выполнять длительный асинхронный сценарий, обращаясь к DOM напрямую, минуя React-рендер на каждом шаге. Через \texttt{forwardRef} и\\ \texttt{useImperativeHandle} экспортируются методы, перечисленные в таблице~\ref{tab:viz_methods}.

\begin{table}[!htbp]
\centering
\scriptsize
\renewcommand{\arraystretch}{1.15}
\begin{tabular}{|p{2.5cm}|p{4.2cm}|p{6.1cm}|}
\hline
\rowcolor{thbg} \textbf{Компонент} & \textbf{Метод} & \textbf{Назначение} \\
\hline
\multirow{8}{*}{\texttt{Visualization}}
& \texttt{show}, \texttt{hide} & показ и скрытие секции \\ \cline{2-3}
& \texttt{getSpeed} & текущая скорость анимации (мс/шаг) \\ \cline{2-3}
& \texttt{setStatus}, \texttt{setOpLabel} & заголовок и метка текущей операции \\ \cline{2-3}
& \texttt{setContainerHTML}, \texttt{appendHTML} & полная или инкрементальная замена содержимого таблицы \\ \cline{2-3}
& \texttt{querySelector}, \texttt{querySelectorAll} & выборка элементов в области таблицы \\ \cline{2-3}
& \texttt{enableBacksub}, \texttt{disableBacksub} & регистрация и отключение второй секции (обратная подстановка) \\ \cline{2-3}
& \texttt{setBacksubStatus}, \texttt{setBacksubContainerHTML}, \texttt{backsubQuery\-Selector} и др. & управление содержимым секции обратной подстановки \\
\hline
\multirow{3}{*}{\texttt{StepLog}}
& \texttt{clear} & очистка журнала \\ \cline{2-3}
& \texttt{addStep} & добавление скрытого шага с возвратом ссылки на DOM-элемент \\ \cline{2-3}
& \texttt{showStep} & асинхронный показ шага через паузу \texttt{stepDelay} (\texttt{Promise}) \\
\hline
\multirow{3}{*}{\texttt{MatrixInput}}
& \texttt{buildGrid} & построение сетки \texttt{<input>}-полей размером $n \times n$ \\ \cline{2-3}
& \texttt{readInput} & чтение данных с валидацией; результат \texttt{\{ n, A, b, extra \}} либо \texttt{null} \\ \cline{2-3}
& \texttt{loadExample} & подстановка демонстрационных данных из контракта \\
\hline
\end{tabular}
\caption{Императивные методы компонентов слоя визуализации}
\label{tab:viz_methods}
\end{table}

    \quad Скорость анимации задаётся пользователем в \texttt{Visualization} выпадающим списком со значениями 3200, 1600 и 800~мс на шаг (соответственно \texttt{0.25x}, \texttt{0.5x}, \texttt{1x}) при значении по умолчанию 1600~мс. Атомарные операции над таблицей собраны в \texttt{solverUtils.js}: \texttt{renderImat} перерисовывает расширенную матрицу как HTML-таблицу с атрибутами \texttt{data-row}/\texttt{data-col}, \texttt{highlightCell} и \texttt{highlightRow} добавляют CSS-классы выделения, \texttt{updateCell} обновляет ячейку с учётом анимационной паузы, а \texttt{clearHighlights} снимает все выделения. Все эти функции обращаются к DOM через \texttt{querySelector} объекта \texttt{viz}, что позволяет повторно использовать их и в основной таблице, и в таблице обратной подстановки.

    \quad В компоненте \texttt{StepLog} метод \texttt{addStep} создаёт скрытый \texttt{<div>} с заголовком, описанием, фрагментом таблицы и связанной BLAS-командой, а \texttt{showStep} ожидает \texttt{stepDelay} миллисекунд, снимает класс \texttt{step--hidden} и прокручивает журнал к новому шагу. Поскольку \texttt{showStep} возвращает \texttt{Promise}, темп журнала контролируется на стороне солвера через \texttt{await}.

    \quad Максимальный порядок задачи зависит от устройства. Тип устройства определяется хуком \texttt{useViewport} (пороги ширины окна 600 и 1024 пикселей), а соответствующие лимиты возвращают функции \texttt{maxMatrixSize},\\ \texttt{maxNonlinearSize} и \texttt{maxRegressionFeatures} (таблица~\ref{tab:adaptive_limits}).

\begin{table}[!htbp]
\centering
\scriptsize
\renewcommand{\arraystretch}{1.15}
\begin{tabular}{|p{5.5cm}|c|c|c|}
\hline
\rowcolor{thbg} \textbf{Класс задачи} & \textbf{Mobile} & \textbf{Tablet} & \textbf{Desktop} \\
\rowcolor{thbg} & \textbf{($\le600$~px)} & \textbf{($\le1024$~px)} & \textbf{(иначе)} \\
\hline
Линейные СЛАУ (порядок) & 5 & 8 & 11 \\
\hline
Нелинейные системы (порядок) & 4 & 7 & 10 \\
\hline
Регрессия (число признаков) & 3 & 4 & 5 \\
\hline
\end{tabular}
\caption{Адаптивные ограничения на размерность задач}
\label{tab:adaptive_limits}
\end{table}

    \quad Дополнительные средства отрисовки нелинейных задач (HTML-блоки якобиана, итерационные таблицы, прогрессбар по параметру гомотопии и т.\,п.) собраны в \texttt{src/utils/nlsRenderHelpers.js}.

\section{Страницы-контейнеры}

    \quad Каждая из трёх страниц-контейнеров реализует единый сценарий взаимодействия (см.\ разд.~2.3): чтение входных данных $\to$ валидация $\to$ проверка \texttt{wasmReady} $\to$ формирование контекста $\to$ \texttt{config.solve(ctx)} $\to$ отображение результата. Состав контекста, передаваемого солверу, приведён в таблице~\ref{tab:ctx_fields}.

\begin{table}[!htbp]
\centering
\scriptsize
\renewcommand{\arraystretch}{1.15}
\begin{tabular}{|p{2.0cm}|p{4.2cm}|p{6.5cm}|}
\hline
\rowcolor{thbg} \textbf{Поле} & \textbf{Источник} & \textbf{Назначение} \\
\hline
\texttt{data} & результат \texttt{readInput} & входные данные задачи в логическом представлении \\ \hline
\texttt{runBlas} & хук \texttt{useWasm} & вызов вычислительного ядра \\ \hline
\texttt{viz} & \texttt{ref} компонента \texttt{Visualization} & императивный интерфейс визуализации \\ \hline
\texttt{stepLog} & \texttt{ref} компонента \texttt{StepLog} & журнал шагов \\ \hline
\texttt{skipRef} & локальный \texttt{useRef} страницы & разделяемый флаг пропуска анимации \\ \hline
\texttt{wasmReady} & хук \texttt{useWasm} & признак готовности ядра \\ \hline
\end{tabular}
\caption{Состав контекста выполнения \texttt{ctx}}
\label{tab:ctx_fields}
\end{table}

    \quad \texttt{SolverPage} обслуживает линейные методы. На монтировании страница загружает контракт через \texttt{configLoader}, монтирует \texttt{MatrixInput},\\ \texttt{Visualization} и \texttt{StepLog} с пропсами из контракта (\texttt{prefix}, \texttt{defaultSize},\\ \texttt{matrixLabel}, \texttt{extraParams}, \texttt{stepDelay}). Перед вызовом \texttt{config.solve(ctx)} страница читает данные методом \texttt{matrixRef.current.readInput()}, проверяет \texttt{wasmReady}, очищает журнал, отключает секцию обратной подстановки и прокручивает контейнер в видимую область.

    \quad \texttt{NonlinearSolverPage} использует $n$ текстовых полей для уравнений (синтаксис вида \texttt{x1\^{}2 + x2\^{}2 - 4}) и поле начального приближения $x_0$. Парсинг и численное дифференцирование возложены на утилиты \texttt{src/utils/mathCore.js}: функция \texttt{buildFunction} компилирует строку выражения в JS-функцию через \texttt{Function} c предварительной валидацией символов по белому списку и подменой математических идентификаторов (\texttt{sin}, \texttt{cos}, \texttt{tan}, \texttt{exp}, \texttt{log}, \texttt{ln}, \texttt{sqrt}, \texttt{abs}, \texttt{pow}, \texttt{PI}, \texttt{E}) на их \texttt{Math}-аналоги; \texttt{evalF} вычисляет значение системы; \texttt{computeJacobian} строит якобиан методом центральных разностей с шагом $h = 10^{-8}$; \texttt{luSolve} решает СЛАУ с частичным выбором ведущего элемента для коррекции $\Delta x$.

    \quad \texttt{RegressionPage} формирует динамическую таблицу наблюдений со строками вида $(x_1, \ldots, x_m, y)$, причём число признаков $m$ выбирается в пределах таблицы~\ref{tab:adaptive_limits}. Минимальное число точек определяется методом \texttt{config.minPoints(m)}, специфичным для каждого солвера. После решения страница получает от солвера структуру \texttt{charts} с предвычисленными данными и в эффекте \texttt{useEffect} отрисовывает её вызовом \texttt{renderCharts(chartsRef.current, charts)}.

\section{Графический модуль}

    \quad Все диаграммы регрессионных методов отрисованы средствами стандартного 2D Canvas API без сторонних библиотек (\texttt{src/utils/chartRenderer.js}, 567 строк). Помощник \texttt{createCanvas(w, h, colors)} создаёт элемент с учётом \texttt{window.devicePixelRatio}: физический размер canvas увеличивается в \texttt{dpr} раз, после чего контекст немедленно масштабируется обратно вызовом\\ \texttt{ctx.scale(dpr, dpr)}, что гарантирует чёткость линий и текста на дисплеях высокой плотности. Функция \texttt{getColors} читает атрибут \texttt{data-theme} корневого элемента документа и возвращает палитру для светлой или тёмной темы, благодаря чему диаграммы автоматически адаптируются при её смене.

    \quad Точка входа \texttt{renderCharts(container, charts)} получает контейнер и объект-описание набора графиков; для каждого присутствующего поля\\ \texttt{charts.*} вызывается соответствующая специализированная функция отрисовки. Список поддерживаемых типов приведён в таблице~\ref{tab:chart_types}.

\begin{table}[H]
\centering
\scriptsize
\renewcommand{\arraystretch}{1.15}
\begin{tabular}{|p{3.6cm}|p{3.7cm}|p{5.4cm}|}
\hline
\rowcolor{thbg} \textbf{Тип} & \textbf{Функция} & \textbf{Применение} \\
\hline
Карточка с формулой модели & \texttt{renderEquationCard} & все методы регрессии \\ \hline
Таблица числовых метрик & \texttt{renderMetrics} & все методы регрессии \\ \hline
Рассеяние с моделью & \texttt{renderScatterChart} & одномерные регрессии ($m=1$) \\ \hline
Частная зависимость & \texttt{renderPartialDependence} & многомерные регрессии ($m \ge 2$) \\ \hline
Предсказание против факта & \texttt{renderPredVsActual} & все методы регрессии \\ \hline
Остатки (рассеяние) & \texttt{renderResidualsScatter} & все методы регрессии \\ \hline
Гистограмма остатков & \texttt{renderResidualHistogram} & все методы регрессии \\ \hline
Коэффициенты модели & \texttt{renderCoefficientsChart} & линейная и полиномиальная регрессии \\ \hline
Тепловая карта корреляций & \texttt{renderCorrelationHeatmap} & многомерные регрессии \\ \hline
\end{tabular}
\caption{Типы диаграмм графического модуля}
\label{tab:chart_types}
\end{table}

    \quad Разделение модельной и графической частей результата (см.\ разд.~2.4) обеспечивается тем, что регрессионный солвер возвращает из \texttt{solve} единый объект, содержащий и параметры модели (коэффициенты, узловые значения, предсказания, остатки, метрики), и набор готовых к отрисовке структур. Благодаря этому одно вычисление обслуживает все диаграммы без повторного обращения к ядру, а смена состава визуализации сводится к изменению набора полей объекта \texttt{charts} и не затрагивает алгоритм.

\chapter{Эксперимент и анализ производительности}

    \quad Глава посвящена количественной оценке реализованной платформы. Эксперимент построен вокруг семи независимых измерений, каждое из которых отвечает на конкретный вопрос о пригодности WebAssembly как среды исполнения численных алгоритмов в браузере: задержка холодной инициализации (§4.2), накладные расходы единичного вызова ядра (§4.3), производительность WebAssembly относительно чистого JavaScript на эквивалентных алгоритмах (§4.4), масштабирование времени решения с порядком задачи (§4.5), эффективность многопоточного режима OpenBLAS (§4.6), стоимость точной рациональной арифметики (§4.7) и накладные расходы слоя визуализации (§4.8). Сводное обсуждение результатов и сопоставление с известными в литературе оценками вынесены в §4.9.

\section{Среда и методика измерений}
\label{sec:exp_method}

    \quad Эксперимент проводится на трёх классах устройств, типичных для целевого окружения платформы (таблица~\ref{tab:exp_testbed}): персональный компьютер на macOS и два мобильных устройства --- iOS и Android. Подобная конфигурация согласуется со схемой исследования энергопотребления WebAssembly в работе~[5], в которой ключевые выводы получены сопоставлением десктопного и мобильного исполнения, а также с измерительной методикой~[4], где сравнение JS и WebAssembly выполнено на широком спектре устройств.

\begin{table}[htbp]
\centering
\scriptsize
\renewcommand{\arraystretch}{1.2}
\begin{tabular}{|p{2.6cm}|p{4.0cm}|p{2.7cm}|p{3.0cm}|}
\hline
\rowcolor{thbg} \textbf{Устройство} & \textbf{Платформа} & \textbf{Браузеры} & \textbf{Многопоточность} \\
\hline
Ноутбук (Mac) & macOS, CPU/RAM/cores --- \textit{TODO} & Chrome, Firefox, Safari & да (4 потока) \\
\hline
Смартфон (Android) & Android \textit{TODO}, ARM \textit{TODO} & Chrome (мобильный) & да (если поддерживается) \\
\hline
Смартфон (iOS) & iOS \textit{TODO}, Apple Silicon & Safari (мобильный) & ограниченная \\
\hline
\end{tabular}
\caption{Состав экспериментального стенда}
\label{tab:exp_testbed}
\end{table}

    \quad \textbf{Инструментарий измерений.} Хронометраж выполняется браузерным API \texttt{performance.now()} (разрешение --- сотни наносекунд при включённой изоляции страницы). Этапы загрузки фиксируются через \texttt{Performance\-Observer} и \texttt{PerformanceNavigationTiming}, что позволяет разложить интервал от \texttt{navigation\-Start} до \texttt{onRuntime\-Initialized} на стадии: загрузка HTML, парсинг JS, скачивание \texttt{shell\_cblas.wasm}, компиляция модуля, инициализация рантайма, готовность пула \texttt{pthread}. Размер бинарника ядра фиксирован (466~КБ, см. §3.1) и одинаков для всех конфигураций.

    \quad \textbf{Контроль повторяемости.} Каждое измерение повторяется $N$ раз ($N \ge 30$ для микробенчмарков и $N \ge 10$ для полных решений СЛАУ); сообщается медиана и межквартильный размах (IQR). Перед серией выполняется один прогон-разогрев, отбрасываемый при анализе. Для измерения холодной загрузки кэш браузера предварительно очищается (включая \texttt{Cache Storage} и память WebAssembly-модуля); для измерения операций после загрузки --- наоборот, прогон выполняется на полностью прогретом V8/JavaScriptCore с целью изолировать издержки самого ядра. Дискретизация по типу устройства реализуется в верстке адаптивно (см. §3.5), но в эксперименте размерности фиксируются явно, а пороги по ширине окна не задействуются.

    \quad \textbf{Ограничения.} Все измерения проводятся над одной и той же сборкой \texttt{shell\_cblas.wasm}, скомпилированной с параметрами из таблицы~\ref{tab:emcc_flags}. Многопоточность OpenBLAS требует включения заголовков \texttt{Cross-Origin-Opener-Policy} и \texttt{Cross-Origin-Embedder-Policy} (для доступа к \texttt{SharedArrayBuffer}); их отсутствие на iOS-Safari в обычных условиях обработано фолбэком на однопоточный режим, что специально отмечено при анализе. Для всех экспериментов источник тестовых данных детерминирован: матрицы и векторы генерируются псевдослучайным генератором с фиксированным seed.

\section{Задержка инициализации платформы}
\label{sec:exp_coldload}

    \quad Эксперимент~1 фиксирует время от первой навигации к платформе до первого момента, когда становится возможным вызов \texttt{runBlas} (флаг \texttt{wasmReady}, см. §3.2). Эта величина характеризует пользовательский опыт и сопоставима с метрикой \emph{ready time} ранее изученных браузерных IDE~[11] и обобщённой оценкой запуска WebAssembly-приложений~[10]. Согласно~[1], типичное время компиляции WebAssembly-модуля порядка~466~КБ должно укладываться в десятки миллисекунд, однако совокупная задержка включает также сетевую загрузку, разбор JS-обёртки и подъём пула pthread.

    \quad В каждой пробе фиксируются пять отметок времени: \texttt{T}$_\text{nav}$ (момент навигации), \texttt{T}$_\text{html}$ (\texttt{DOMContent\-Loaded}), \texttt{T}$_\text{wasm-fetch}$ (завершение скачивания \texttt{shell\_cblas.wasm}), \texttt{T}$_\text{init}$ (вызов \texttt{onRuntime\-Initialized}) и \texttt{T}$_\text{ready}$ (выставление флага \texttt{wasmReady} в \texttt{WasmContext}). По разностям этих отметок эксперимент выделяет четыре стадии: \emph{скачивание}, \emph{компиляция}, \emph{инициализация рантайма}, \emph{подъём пула потоков}. Измерение проводится отдельно для холодного (свежий запуск, очищенный кэш) и тёплого (повторное открытие) сценариев. Результат представлен на рис.~\ref{fig:exp_coldload}.

\begin{figure}[htbp]
\centering
\begin{tikzpicture}[font=\footnotesize]
  \draw[gray!40, very thin] (0,0) grid[xstep=2,ystep=1] (12,5);
  \draw[thick, ->] (0,0) -- (12.3,0) node[right]{браузер $\times$ кэш};
  \draw[thick, ->] (0,0) -- (0,5.3) node[above]{$T$, мс};
  \foreach \x/\l in {1/Chrome cold, 3/Chrome warm, 5/Firefox cold, 7/Firefox warm, 9/Safari cold, 11/Safari warm}
     \node[below, rotate=20, anchor=east, xshift=3pt] at (\x,-0.1) {\l};
  \node[align=center, text=gray] at (6,2.5) {\itshape (стопка из 4 стадий: скачивание, компиляция,\\инициализация рантайма, пул потоков; данные --- TODO)};
\end{tikzpicture}
\caption{Задержка инициализации платформы по стадиям, разбивка по браузерам и сценариям кэша}
\label{fig:exp_coldload}
\end{figure}

\section{Накладные расходы вызова ядра}
\label{sec:exp_boundary}

    \quad Эксперимент~2 измеряет фиксированную составляющую задержки одного вызова \texttt{runBlas} --- последовательность из шести шагов (см. таблицу~\ref{tab:runblas_pipeline}), исполняемая независимо от объёма передаваемых данных. Эта величина критична для алгоритмов, в которых вычисление складывается из большого числа мелких операций; в работе~[3] показано, что именно граничные издержки JS$\leftrightarrow$WASM являются основным источником разрыва WebAssembly с нативной производительностью на коротких задачах, а в~[7] подчёркнута важность их учёта при интеграции WebAssembly в производственные веб-приложения.

    \quad В качестве пробной операции выбрана \texttt{idamax} над вектором длины $n$, варьируемой по геометрической прогрессии ($n = 10, 50, 100, 500, 1000, 5000$). На каждом значении $n$ выполняется $K = 1000$ повторов и сохраняется медианное время одного вызова. Для отделения собственно граничной составляющей от полезной работы дополнительно измеряется чисто JavaScript-эквивалент: проход по массиву поиском индекса максимального по модулю элемента. Разница между медианами при $n \to 0$ оценивает ``пол'' задержки --- сумму расходов на \texttt{lengthBytesUTF8}, \texttt{\_malloc}, \texttt{stringToUTF8}, \texttt{\_run\_command}, \texttt{\_free} и перехват вывода. Результат --- рис.~\ref{fig:exp_boundary}.

\begin{figure}[htbp]
\centering
\begin{tikzpicture}[font=\footnotesize]
  \draw[gray!40, very thin] (0,0) grid[xstep=2,ystep=1] (12,5);
  \draw[thick, ->] (0,0) -- (12.3,0) node[right]{$\log_{10} n$};
  \draw[thick, ->] (0,0) -- (0,5.3) node[above]{$T$/вызов, мкс};
  \foreach \x/\l in {1/1, 3/2, 5/2.7, 7/3, 9/3.7, 11/4}
     \node[below] at (\x,0) {\l};
  \node[align=center, text=gray] at (6,2.5) {\itshape (две кривые: \texttt{runBlas(idamax)} и чистый JS;\\разрыв при малых $n$ --- ``пол'' граничных издержек; данные --- TODO)};
\end{tikzpicture}
\caption{Время одного вызова \texttt{runBlas(idamax)} в сравнении с эквивалентом на чистом JavaScript}
\label{fig:exp_boundary}
\end{figure}

\section{Сравнение WebAssembly и JavaScript на одном алгоритме}
\label{sec:exp_wasmjs}

    \quad Эксперимент~3 сопоставляет производительность ядра и чистого JavaScript на алгоритмически эквивалентной задаче --- решении системы $Ax = b$ методом исключения Гаусса с частичным выбором ведущего элемента. Аналогичные сравнения для матрично-векторных операций приведены в~[4] и~[6], причём оба источника отмечают существенное преимущество WebAssembly над оптимизированным JS на матричных задачах среднего и большого порядка; в~[10] поставлен вопрос о том, насколько это преимущество устойчиво в реальных веб-приложениях.

    \quad WebAssembly-сторону представляет команда \texttt{dgesv} (полное решение СЛАУ через LAPACK). JavaScript-сторону --- функция \texttt{luSolve} из \texttt{src/utils/mathCore.js} (LU-факторизация с частичным выбором ведущего, в коде проекта используется в нелинейных солверах --- см. §3.6). Размерности: $n \in \{10, 50, 100, 200, 500, 1000\}$. На каждом $n$ генерируется $M = 10$ случайных хорошо обусловленных матриц с фиксированным seed; время решения каждой матрицы измеряется отдельно. На рис.~\ref{fig:exp_wasmjs_abs} приведено абсолютное время решения, на рис.~\ref{fig:exp_wasmjs_speedup} --- ускорение $T_\text{JS} / T_\text{WASM}$ как функция $n$.

\begin{figure}[htbp]
\centering
\begin{tikzpicture}[font=\footnotesize]
  \draw[gray!40, very thin] (0,0) grid[xstep=2,ystep=1] (12,5);
  \draw[thick, ->] (0,0) -- (12.3,0) node[right]{$\log_{10} n$};
  \draw[thick, ->] (0,0) -- (0,5.3) node[above]{$\log_{10} T$, мс};
  \foreach \x/\l in {1/1, 3/1.7, 5/2, 7/2.3, 9/2.7, 11/3}
     \node[below] at (\x,0) {\l};
  \node[align=center, text=gray] at (6,2.5) {\itshape (две кривые: \texttt{dgesv} (WASM) и \texttt{luSolve} (JS),\\log--log; данные --- TODO)};
\end{tikzpicture}
\caption{Время решения СЛАУ при варьировании порядка матрицы}
\label{fig:exp_wasmjs_abs}
\end{figure}

\begin{figure}[htbp]
\centering
\begin{tikzpicture}[font=\footnotesize]
  \draw[gray!40, very thin] (0,0) grid[xstep=2,ystep=1] (12,5);
  \draw[thick, ->] (0,0) -- (12.3,0) node[right]{$n$};
  \draw[thick, ->] (0,0) -- (0,5.3) node[above]{$T_\text{JS}/T_\text{WASM}$};
  \draw[dashed, gray] (0,1) -- (12,1) node[right, black]{$=1$};
  \foreach \x/\l in {1/10, 3/50, 5/100, 7/200, 9/500, 11/1000}
     \node[below] at (\x,0) {\l};
  \node[align=center, text=gray] at (6,3) {\itshape (одна кривая ускорения; разрыв при малых $n$,\\рост при больших; данные --- TODO)};
\end{tikzpicture}
\caption{Ускорение WebAssembly относительно JavaScript на эквивалентном алгоритме}
\label{fig:exp_wasmjs_speedup}
\end{figure}

\section{Масштабирование времени решения с порядком задачи}
\label{sec:exp_scaling}

    \quad Эксперимент~4 верифицирует асимптотическую сложность реализованных в платформе классов методов: прямого (Гаусс, $O(n^3)$) и итерационных (Якоби, GMRES, $O(n^2 \cdot k)$, где $k$ --- число итераций до сходимости). Цель --- проверить, что наблюдаемый показатель степени $\alpha$ в законе $T(n) \propto n^\alpha$ согласуется с теоретической оценкой и что переход на WebAssembly не вносит дополнительных скрытых членов сложности; постановка близка по духу к~[6]. Дополнительная мотивация --- получение эмпирических ориентиров для пользователя при выборе метода в зависимости от размерности задачи.

    \quad Размерности $n \in \{10, 50, 100, 200, 500, 1000\}$ (Гаусс) и $n \in \{50, 100, 500, 1000\}$ (GMRES) при условии хорошей обусловленности. Время измеряется аналогично §4.4. По результатам каждого метода строится регрессия $\log T = \alpha \log n + \beta$, и оцениваемое значение $\alpha$ сравнивается с теоретическим. Результат --- рис.~\ref{fig:exp_scaling}.

\begin{figure}[htbp]
\centering
\begin{tikzpicture}[font=\footnotesize]
  \draw[gray!40, very thin] (0,0) grid[xstep=2,ystep=1] (12,5);
  \draw[thick, ->] (0,0) -- (12.3,0) node[right]{$\log_{10} n$};
  \draw[thick, ->] (0,0) -- (0,5.3) node[above]{$\log_{10} T$, мс};
  \node[align=center, text=gray] at (6,2.5) {\itshape (две--три кривые: Гаусс, Якоби, GMRES;\\наклон оценивает показатель сложности; данные --- TODO)};
\end{tikzpicture}
\caption{Зависимость времени решения от порядка матрицы для прямого и итерационных методов}
\label{fig:exp_scaling}
\end{figure}

\section{Эффективность многопоточности}
\label{sec:exp_threads}

    \quad Эксперимент~5 оценивает выигрыш от параллельного исполнения OpenBLAS на ядре. Поддержка многопоточности заявлена через параметры сборки \texttt{USE\_THREAD=1}, \texttt{NUM\_THREADS=4}, \texttt{PTHREAD\_POOL\_SIZE=4} (см. таблицу~\ref{tab:emcc_flags}); число активных потоков переключается командой ядра \texttt{threads $n$}. Постановка опирается на~[1], где многопоточная модель WebAssembly обоснована как ключевой фактор приближения к нативной производительности на ресурсоёмких вычислениях, и на~[5], где энергоэффективность обсуждается с учётом параллелизма.

    \quad В качестве ресурсоёмкой операции выбрана команда \texttt{dgemm} над квадратными матрицами порядка $500$. Перебираются конфигурации $n_\text{th} \in \{1, 2, 4\}$ при прочих равных. На мобильных устройствах фиксируется фактическая поддержка многопоточности браузером (наличие \texttt{SharedArrayBuffer} и корректных заголовков COOP/COEP); при недоступности измерения проводятся в однопоточном режиме и эта особенность отмечается при анализе. Метрика --- ускорение $S(n_\text{th}) = T(1) / T(n_\text{th})$. Результат --- рис.~\ref{fig:exp_threads}.

\begin{figure}[htbp]
\centering
\begin{tikzpicture}[font=\footnotesize]
  \draw[gray!40, very thin] (0,0) grid[xstep=2,ystep=1] (10,5);
  \draw[thick, ->] (0,0) -- (10.3,0) node[right]{число потоков $n_\text{th}$};
  \draw[thick, ->] (0,0) -- (0,5.3) node[above]{$S(n_\text{th})$};
  \draw[dashed, gray] (0,0) -- (8,4) node[right, black]{идеальное};
  \foreach \x/\l in {2/1, 5/2, 8/4}
     \node[below] at (\x,0) {\l};
  \node[align=center, text=gray] at (5,4) {\itshape (столбцы для Mac/Android/iOS;\\идеальная диагональ для сравнения; данные --- TODO)};
\end{tikzpicture}
\caption{Ускорение операции \texttt{dgemm} ($n=500$) в зависимости от числа потоков OpenBLAS}
\label{fig:exp_threads}
\end{figure}

\section{Стоимость точной арифметики}
\label{sec:exp_exact}

    \quad Эксперимент~6 оценивает накладные расходы режима точной арифметики (см. §3.4). Поскольку задача поддержки рациональных дробей решается силами JavaScript, цена точности слагается из двух частей: стоимость обращений к классу \texttt{Frac} вместо нативных арифметических операций IEEE~754 и невозможность делегировать шаги алгоритма ядру. В литературе эта величина не имеет прямого аналога, поскольку точная арифметика в JS-средах обычно реализуется отдельными библиотеками без интеграции с пошаговой визуализацией; настоящий опыт --- оригинальный вклад работы.

    \quad Метод~--- Гаусс, режим выбирается предикатом \texttt{isAllInt} на этапе чтения данных. Размерности $n \in \{3, 5, 7, 9, 11\}$ (верхняя граница соответствует адаптивному лимиту настольной версии, таблица~\ref{tab:adaptive_limits}). На каждом $n$ генерируется $M = 20$ случайных матриц с целыми элементами в диапазоне $[-9, 9]$ и проверяется выполнение точного равенства $A x = b$ в поле $\mathbb{Q}$. Сравниваются медианы $T_\text{exact}(n)$ и $T_\text{float}(n)$, а также абсолютная ошибка решения относительно эталонного, полученного через \texttt{dgesv}. Результат --- рис.~\ref{fig:exp_exact}.

\begin{figure}[htbp]
\centering
\begin{tikzpicture}[font=\footnotesize]
  \draw[gray!40, very thin] (0,0) grid[xstep=2,ystep=1] (10,5);
  \draw[thick, ->] (0,0) -- (10.3,0) node[right]{$n$};
  \draw[thick, ->] (0,0) -- (0,5.3) node[above]{$T_\text{exact}/T_\text{float}$};
  \foreach \x/\l in {2/3, 4/5, 6/7, 8/9, 10/11}
     \node[below] at (\x,0) {\l};
  \node[align=center, text=gray] at (5,2.5) {\itshape (отношение времени точного режима к численному;\\рост при увеличении $n$; данные --- TODO)};
\end{tikzpicture}
\caption{Относительная стоимость режима точной арифметики (\texttt{Frac}) против численного режима}
\label{fig:exp_exact}
\end{figure}

\section{Издержки слоя визуализации}
\label{sec:exp_viz}

    \quad Эксперимент~7 разделяет полное время отклика метода на две составляющие: чистое вычисление и накладные расходы пошаговой визуализации (анимационные паузы, обновление DOM, журнал шагов). Постановка близка по идее к~[8] и~[9], где обсуждается стоимость интерактивной визуализации в браузере для научных приложений; ключевое отличие данной работы --- зависимость общего времени от пользовательской настройки скорости анимации (3200/1600/800~мс на шаг, см. §3.5).

    \quad Сравниваются три сценария на одном и том же методе (Гаусс, $n = 8$, фиксированные данные): (а) \emph{чистое вычисление} --- выполнение алгоритма без обращений к \texttt{viz} и \texttt{stepLog}; (б) \emph{визуализация с пропуском} --- вызов \texttt{handleSkip} немедленно после старта (\texttt{skipRef.current = true}); (в) \emph{визуализация на максимальной скорости} (800~мс на шаг). Время измеряется по \texttt{performance.now()} от начала \texttt{config.solve(ctx)} до завершения промиса. Результат --- рис.~\ref{fig:exp_viz}.

\begin{figure}[htbp]
\centering
\begin{tikzpicture}[font=\footnotesize]
  \draw[gray!40, very thin] (0,0) grid[xstep=2,ystep=1] (10,5);
  \draw[thick, ->] (0,0) -- (10.3,0) node[right]{сценарий};
  \draw[thick, ->] (0,0) -- (0,5.3) node[above]{$T$, мс};
  \foreach \x/\l in {2/чистое, 5/{с пропуском}, 8/{полная анимация}}
     \node[below, align=center] at (\x,0) {\l};
  \node[align=center, text=gray] at (5,2.5) {\itshape (стопочные столбцы: компонент компьютер.,\\анимационные паузы, DOM-обновления; данные --- TODO)};
\end{tikzpicture}
\caption{Разложение полного времени отклика метода на компоненты}
\label{fig:exp_viz}
\end{figure}

\section{Обсуждение и сопоставление с известными результатами}
\label{sec:exp_discussion}

    \quad Полная интерпретация будет приведена после получения количественных данных. Ниже зафиксированы ожидаемые качественные выводы и точки сопоставления с литературой.

    \quad \textbf{Загрузка и инициализация (§4.2).} Согласно~[1], характерное время компиляции WebAssembly-модуля размером порядка сотен килобайт составляет десятки миллисекунд; основной вклад в полную задержку ожидается со стороны сетевого скачивания и подъёма пула потоков. На тёплом кэше следует ожидать сокращения времени до десятков миллисекунд, что соответствует наблюдениям~[11] для браузерных IDE.

    \quad \textbf{Граничные издержки (§4.3).} В работе~[3] установлено, что фиксированный пол задержки одного вызова через границу JS$\leftrightarrow$WASM является доминирующим фактором замедления для коротких операций. Для текстового командного протокола настоящей работы пол ожидаемо выше, чем для прямого FFI, за счёт расходов на кодирование UTF-8 и аллокацию буфера; критическая размерность, при которой обращение к ядру становится оправданным, --- предмет численной оценки в §4.3.

    \quad \textbf{WebAssembly против JavaScript (§4.4--§4.5).} Результаты~[4] и~[6] предсказывают существенное ($> 2\times$) преимущество WebAssembly при $n \gtrsim 100$. Настоящий эксперимент позволит локализовать пересечение, при котором WebAssembly начинает выигрывать у JavaScript-реализации с учётом граничных издержек, и оценить, насколько это значение согласуется с указанными работами на конкретной задаче решения СЛАУ.

    \quad \textbf{Многопоточность (§4.6).} Многопоточная модель WebAssembly~[1] на персональных компьютерах ожидаемо даёт ускорение, близкое к идеальному при числе потоков, не превышающем число физических ядер. На мобильных устройствах ожидается частичный регресс из-за тепловых ограничений и неполной поддержки \texttt{SharedArrayBuffer}, что согласуется с энергетическим анализом~[5].

    \quad \textbf{Точная арифметика (§4.7).} Поскольку точный режим не использует ядро, увеличение времени должно следовать кубическому росту числа операций над \texttt{Frac} с дополнительным фактором роста знаменателей. Этот эксперимент --- оригинальный вклад работы; полученные данные позволят уточнить рекомендуемые границы применимости точного режима в учебных сценариях.

    \quad \textbf{Визуализация (§4.8).} Доля анимационных пауз и DOM-обновлений в полном времени отклика характеризует пригодность платформы как интерактивного учебного инструмента; ожидается, что эта доля доминирует на малых $n$, где собственно вычисление занимает доли миллисекунды, и сокращается при увеличении размера задачи.

    \quad Окончательные выводы и количественное сопоставление с~[1, 3, 4, 5, 6, 10] вынесены в Заключение после заполнения таблиц и графиков настоящей главы.

\chapter*{Заключение}

\renewcommand{\bibname}{\large Список литературы}
\begin{thebibliography}{9}
        \vspace*{0.3cm}
        \bibitem{leung}
        Joseph Y-T. Leung\ Handbook of Scheduling: Algorithms, Models, and
        Performance Analysis CRC Press 2004
        \bibitem{galin} Michelle Galin, Anna Hedblom Linux Kernel
        Scheduler Evaluation
        for Performance-Critical Telecom Workloads Linköping University 2025
        \bibitem{mao} Jinsong Mao, Benjamin E. Ujcich, Shiqing Ma
        Rethinking Provenance
        Completeness with a Learning-Based Linux Scheduler arXiv 2025
        \bibitem{wang} Xinbo Wang, Shian Jia, Ziyang Huang, Jing Cao,
        Mingli Song
        Mixture-of-Schedulers: An Adaptive Scheduling Agent as a
        Learned Router for
        Expert Policies arXiv 2025
        \bibitem{stoica}
        Ion Stoica, Hussein Abdel-Wahab Earliest Eligible Virtual Deadline First
        : A Flexible and Accurate Mechanism for Proportional Share
        Resource Allocation
        Old Dominion University 1996
        \bibitem{shago} A1349 Pavel Shago [Электронный ресурс].
        -- Режим доступа: \url{https://github.com/shgpavel/A1349} --
        Дата обращения
        16.12.2025.
        \bibitem{tanenbaum} Andrew S. Tanenbaum, Herbert Bos\ MODERN OPERATING
        SYSTEMS Pearson
        Education 2023
        \bibitem{rice}
        Liz Rice Learning eBPF O’Reilly Media, Inc 2023
        \bibitem{torvalds}
        Torvalds L. Linux Kernel [Электронный ресурс] / L. Torvalds.
        -- Режим доступа: \url{https://github.com/torvalds/linux} --
        Дата обращения
        04.12.2025.
        \bibitem{scx}
        sched-ext/scx [Электронный ресурс]. -- Режим доступа: \\
        \url{https://github.com/sched-ext/scx} -- Дата обращения 07.12.2025.
        \bibitem{lkml}
        lkml [Электронный ресурс]. -- Режим доступа: \url{https://lkml.org}
        -- Дата обращения 24.11.2025.
        \bibitem{humphries}
        Jack Tigar Humphries, Neel Natu, Ashwin Chaugule, Ofir Weisse,
        Barret Rhoden, Josh Don, Luigi Rizzo, Oleg Rombakh, Paul Turner,
        Christoph Lameter\ ghOST: Scheduling Beyond the OS
        ACM SIGOPS 2021
\end{thebibliography}

\end{document}